\documentclass[11pt,a4paper,oneside]{report}
\usepackage[margin=27mm]{geometry}
\usepackage[T1]{fontenc}
\usepackage[utf8]{inputenc}
\usepackage{lmodern,microtype,amsmath,amssymb,amsthm,bm}
\usepackage{booktabs,longtable,array,graphicx,xcolor,enumitem}
\usepackage[round,authoryear]{natbib}
\usepackage{tikz}
\usetikzlibrary{arrows.meta,positioning}
\usepackage[colorlinks=true,linkcolor=blue!50!black,citecolor=blue!50!black,urlcolor=blue!50!black]{hyperref}
\usepackage{bookmark}
\hypersetup{pdftitle={Cash Flows, Macroeconomic Scenarios, and FX Hedging},pdfauthor={Research proposal}}
\setlength{\parskip}{5pt}
\setlength{\parindent}{0pt}
\setlength{\emergencystretch}{2em}
\setcounter{tocdepth}{1}
\numberwithin{equation}{chapter}
\newcommand{\E}{\mathbb E}
\newcommand{\V}{\operatorname{Var}}
\newcommand{\Cov}{\operatorname{Cov}}
\newcommand{\N}{\mathcal N}
\newcommand{\ind}{\mathbf 1}
\newcommand{\tr}{\operatorname{tr}}
\newcommand{\diag}{\operatorname{diag}}
\newcommand{\KL}{\operatorname{KL}}
\newcommand{\softplus}{\operatorname{softplus}}
\newcommand{\logit}{\operatorname{logit}}
\newcommand{\sigmoid}{\operatorname{sigmoid}}
\newcommand{\Hist}{\mathcal H}
\newcommand{\Avail}{\mathcal A}
\newcommand{\doop}{\operatorname{do}}
\newcommand{\argmax}{\operatorname*{arg\,max}}
\newcommand{\argmin}{\operatorname*{arg\,min}}
\newtheorem{proposition}{Proposition}[chapter]
\begin{document}
\hypersetup{pageanchor=false}
\begin{titlepage}
\vspace*{18mm}
{\Huge\bfseries Cash Flows,\\ Macroeconomic Scenarios,\\ and FX Hedging\par}
\vspace{8mm}
{\Large A literature survey and phased product proposal for commercial banking\par}
\vspace{10mm}
{\large Version 2 \quad 9 October 2026\par}
\vspace{15mm}
\textbf{The decision.} Which clients are likely to face material foreign-currency exposure, which available hedge would improve their financial position, and where can a limited sales team create incremental value?

\textbf{The proposed design.} A shared dynamic model forecasts currency-specific cash receipts and payments under explicit macroeconomic scenarios. An independently estimated utility and discrete-choice model translates those distributions into FX-hedging decisions. A later causal evaluation measures the incremental value of outreach.

\textbf{The evidence.} Twelve primary works are reconstructed in the notation of the banking problem. All cited papers are stored locally. Algebraic checks and document compilation accompany this proposal; bank-data validation and causal identification remain work to be undertaken.
\vfill
Prepared for quantitative research, treasury product, relationship management, and model-risk teams. Assumed prerequisites: probability, calculus, linear algebra, and regression. Numerical examples are illustrative, not estimates for any bank or client.
\end{titlepage}
\pagenumbering{roman}
\hypersetup{pageanchor=true}
\chapter*{Executive proposal}
\addcontentsline{toc}{chapter}{Executive proposal}
A bank with one million business clients and three years of monthly observations has a large panel and a short history. Shared estimation can learn recurring differences between clients, industries, and currencies. It cannot manufacture additional business cycles. The central proposal is therefore a parsimonious, hierarchical state-space model of gross cash inflows and outflows, with a separately specified macroeconomic scenario channel and explicit uncertainty about the portion of each client's activity observed by the bank.

Forecasting receipts and payments separately preserves the financial information needed for hedging. A client expecting EUR 1 million of receipts and EUR 0.8 million of payments at the same date has a different exposure from a client expecting EUR 1 million of receipts alone. Their net exposures depend on timing, existing hedges, balances, and the covariance between receipts and payments. A binary transaction forecast discards these quantities. A joint probability distribution over future flows can instead support exposure measurement, stress analysis, product design, and the eventual allocation of sales capacity.

The macroeconomic input is an explicit scenario path: for example, real GDP growth, rates, inflation, and FX conditions. The requested positive GDP effect on every client's expected inflow can be imposed through nonnegative response coefficients and lag weights. That restriction makes the scenario behavior transparent. It is an economic assumption to test, not a universal empirical law. A conditional forecast becomes a causal counterfactual only after the relevant structural equations, confounding assumptions, and invariance under intervention have been justified.

The FX module consumes forecast paths but does not determine the cash-flow model's training objective. Corporate-finance arguments explain why smoothing internal funds may be valuable when external finance is costly or liquidity shortfalls disrupt investment. A tractable certainty-equivalent calculation then values each feasible hedge, including its price and transaction costs. Random utility yields a discrete-choice model over actual available offers and an outside option. Choice-set records, existing hedge positions, indicative quotes, and declined offers are consequently as important as transaction histories for the second phase.

Delivery proceeds through four useful increments: a reliable currency-flow data set and baselines; joint probabilistic cash-flow forecasts with scenario analysis; FX exposure and client-value calculations; and estimated choice plus randomized outreach evaluation. The first forecasting product can be delivered without waiting for preference identification. Neither model accuracy nor sales uplift is established by this document. The proposed validation compares against simple cash-flow heuristics, checks joint tail behavior and macro sensitivity, and ultimately measures incremental net contribution under real capacity constraints.
\tableofcontents
\clearpage
\pagenumbering{arabic}
\chapter{The updated banking problem}
\label{ch:problem}
\section{What must be predicted, and what must be decided?}
The bank wants to use a limited relationship-management team to reach clients before a financially consequential need arises. A scheduled payment is useful evidence, but the commercial decision concerns the client's exposure and the value of a suitable product. The first task is to forecast financial quantities; the second is to evaluate a hedge; the third is to predict a choice; and the fourth is to decide whether contacting the client changes the outcome enough to justify scarce sales time. These tasks share information but have different targets.

Let $i=1,\ldots,n$ identify legal entities, with $n$ approximately one million. Let $t=1,\ldots,T$ index month ends and $T=36$ denote the available historical length. Currency is indexed by $c$, and the superscripts $+$ and $-$ denote gross receipts and payments. The monthly quantities
\begin{equation}
F^+_{ict}\geq0,\qquad F^-_{ict}\geq0
\end{equation}
are measured in native currency units. They are economic cash flows across the chosen client boundary, before financial hedges, after removing duplicate records and transfers that are internal to that boundary. Their observed portions, $Y^+_{ict}$ and $Y^-_{ict}$, are what pass through the bank. The distinction matters throughout the proposal.

The information set $\Hist_{it}$ contains only information actually available when the decision is made: observed flows to date, dated static and industry attributes $z_i$, past macroeconomic releases, dated sales activity, and verified existing positions. Let $m_{t+1:t+H}$ be an externally supplied macroeconomic path over a horizon $H$. The first deliverable is a joint conditional distribution
\begin{equation}
\mathcal P_{it}^{m}=
 p\left(\{Y^+_{ic,t+h},Y^-_{ic,t+h}\}_{c,h=1}^{C,H}
 \mid\Hist_{it},m_{t+1:t+H}\right).
\label{eq:forecasttarget}
\end{equation}
The notation on the right means every selected currency and every future month. It is a distribution of paths, not just a vector of independent marginal forecasts. If a separate observation model can identify total-wallet flows, $F$ can replace $Y$. Until then, a bank-observed-flow forecast is the estimable target.

Three-month exposure provides a running example. Suppose an exporter expects EUR 1 million in receipts and EUR 0.8 million in supplier payments at the same settlement date. Its unhedged net operating exposure is EUR 0.2 million, before balances and existing hedges. A forward sale of EUR 1 million could substantially overhedge the position. If the supplier payments occur two months earlier, the same totals also conceal a liquidity problem. The joint forecast must retain the currency, sign, and settlement bucket of each flow to support either decision.

\section{From flows to utility, choice, and sales capacity}
Let $j$ denote an available hedge offer with currency, maturity, notional, forward rate, collateral terms, fees, and service attributes. Let $W_{it}(j,\omega)$ be the client's financial outcome under future state $\omega$. A utility specification defines
\begin{equation}
V_{itj}=\mathcal U_i\bigl(W_{it}(j,\cdot);\mathcal P_{it}^{m}\bigr)
           -\text{nonfinancial friction}_{itj}.
\label{eq:valueintro}
\end{equation}
The functional $\mathcal U_i$ may be expected continuation value, a certainty equivalent, or an explicitly approximate mean--variance criterion. It measures client value, in a declared unit. It is not the bank's sales margin. Random utility and the available set $\Avail_{it}$ produce probabilities
\begin{equation}
p_{itj}=\Pr\{j=\argmax_{k\in\Avail_{it}}(V_{itk}+\epsilon_{itk})
                 \mid\Hist_{it},\Avail_{it}\}.
\end{equation}
The outside option may include waiting, accepting exposure, or dealing elsewhere. When those actions are not separately observed, their composition cannot be identified from a single ``no trade at our bank'' label.

Finally, let $A_{it}\in\{0,1\}$ denote an outreach assignment and $R_{it}(a)$ the bank's net contribution under that assignment. A client can have high forecast exposure and high purchase probability but little incremental response to a call. The final commercial target is
\begin{equation}
\tau_{it}=\E[R_{it}(1)-R_{it}(0)\mid\Hist_{it}],
\qquad
\max_{a_i\in\{0,1\}}\sum_i a_i\tau_{it}
\quad\text{subject to}\quad\sum_i d_i a_i\leq B_t,
\label{eq:capacityintro}
\end{equation}
where $d_i$ is expected sales time and $B_t$ is available capacity. Eligibility, client ownership, and suitability constraints can further restrict feasible assignments. This objective explains why forecasting, product choice, and causal outreach effects should remain separate modules.

\begin{figure}[ht]
\centering
\begin{tikzpicture}[node distance=5mm,box/.style={draw,rounded corners,align=center,text width=11.5cm,inner sep=7pt},>=Stealth]
\node[box] (a) {As-of-date observations: receipts, payments, currency,\nobreakspace macro releases, client attributes};
\node[box,below=of a] (b) {Joint cash-flow model\\Shared dynamics and explicit macroeconomic scenario responses};
\node[box,below=of b] (c) {Exposure and utility module\\Settlement buckets, balances, existing hedges, feasible FX offers};
\node[box,below=of c] (d) {Choice module\\Available alternatives, prices, preferences, outside option};
\node[box,below=of d] (e) {Outreach policy\\Incremental bank contribution subject to sales capacity};
\draw[->] (a)--(b);\draw[->] (b)--(c);\draw[->] (c)--(d);\draw[->] (d)--(e);
\end{tikzpicture}
\caption{Each stage answers a distinct question. Scenario paths and uncertainty pass forward; the forecasting stage can be delivered before preferences or outreach effects are estimated.}
\end{figure}

\section{Why a million clients do not eliminate the short-history problem}
The cross-section helps estimate shared seasonality, industry differences, and relationships between prior activity and future flows. It does not supply a million independent realizations of GDP. To see the limitation, consider
\begin{equation}
y_{it}=\alpha_i+\beta g_t+u_t+e_{it},
\qquad \E[e_{it}\mid g,u]=0,
\end{equation}
where $u_t$ is an unobserved common demand shock. Averaging across clients gives
\begin{equation}
\overline y_t=\overline\alpha+\beta g_t+u_t+\overline e_t.
\end{equation}
With weak cross-sectional dependence, $\overline e_t$ can shrink as $n$ grows. The common shock $u_t$ remains, as do only 36 values of $g_t$. If $g_t$ and $u_t$ are correlated, cross-sectional averaging does not remove bias. If unrestricted month effects are included, they span $g_t$ exactly and prevent separate identification of a common GDP coefficient. The proposed model must therefore state which common movements are attributed to GDP and which to residual factors.

Three annual cycles also provide little evidence about stable client-specific seasonality, and potentially no evidence about a recession or a sharp exchange-rate discontinuity. Hierarchical pooling is a response to this information constraint. It should be combined with simple baselines, uncertainty intervals, and sensitivity to macroeconomic assumptions.

\section{What the literature can establish}
The survey connects four bodies of work. Global forecasting explains parameter sharing across short series. State-space estimation explains how noisy observations update latent cash-flow conditions. Structural causal models explain the assumptions that distinguish scenarios from interventions. Corporate hedging and discrete-choice theory connect exposures to economically meaningful product decisions. Policy learning supplies the final distinction between a likely customer and a customer whose outcome can be changed by outreach.

The cited methods are derived at the level used in the proposal, including likelihoods, inference recursions, utility calculations, and choice estimators. The aim is a self-contained reconstruction of their relevant methods, rather than a reproduction of every theorem and empirical application in each source. Source-specific technical anchors and version differences appear alongside the derivations and in Chapter~\ref{ch:sources}. The resulting architecture is a proposed synthesis. No cited study establishes its predictive or commercial performance on this bank's data.

\chapter{Cash-flow measurement and the information the bank lacks}
\label{ch:observation}
\section{Define a financial boundary before fitting a model}
For an individual account, a transfer to another account of the same company is an outflow. For the consolidated legal entity it may be internal. For a multinational group it may be an intercompany funding movement with different tax and treasury implications. The project must select the entity boundary, retain account-to-entity mappings, and flag uncertain links. Consolidating by similar names alone is inadequate: false merges can create fictitious natural hedges between unrelated firms.

Native-currency gross flows should distinguish operating receipts and payments, financing, asset purchases, internal transfers, FX conversions, and settlements of existing derivatives when the bank can identify them. A sale of euros creates two currency legs; counting both as new operating activity would double count the economic event. Removing derivative settlements from the operating-flow target also prevents the model from treating its own hedge recommendations as new commercial demand.

At month end, a reconciliation for a currency account is
\begin{equation}
B_{ic,t}-B_{ic,t-1}=Y^+_{ict}-Y^-_{ict}+J_{ict},
\label{eq:accounting}
\end{equation}
where $J$ contains explicitly identified adjustments, interest, reclassifications, and any transactions excluded by the selected target definition. This is an accounting check, not a behavioral equation. FX translation changes the USD reporting value of a EUR balance without changing the number of euros. Modeling only USD-converted transaction amounts confounds volume with exchange rates.

Monthly aggregation limits the product. Two offsetting payments within a month can produce zero monthly net flow while creating a large intramonth liquidity need. With only monthly data, the first release can support monthly exposure windows and conversations about future needs. A claim to optimize daily hedge execution or intramonth liquidity would require finer settlement information.

\section{A simple proof of the total-wallet identification problem}
Suppose for the moment that amounts observed at the bank satisfy
\begin{equation}
Y^+_{ict}=s^+_{ict}F^+_{ict},\qquad
Y^-_{ict}=s^-_{ict}F^-_{ict},\qquad 0<s^\pm_{ict}\leq1.
\label{eq:capture}
\end{equation}
These equations are an illustrative observation model, not an assertion that a bank receives a constant fraction of every transaction. They already show the problem. For any $a\geq1$, replacing $(F,s)$ by $(aF,s/a)$ leaves $Y$ unchanged. No estimator using $Y$ alone can distinguish these alternatives. In logs, the same result is
\begin{equation}
\log Y=\log F+\log s;
\end{equation}
a change in total economic activity and a change in bank share can be observationally identical.

Separate capture shares for inflows and outflows are especially consequential. Observed receipts of EUR 1 million and payments of EUR 0.8 million imply a bank-observed net EUR 0.2 million. If $s^+=0.5$ and $s^-=0.8$, total receipts and payments are EUR 2 million and EUR 1 million, so total net exposure is EUR 1 million. If the shares are reversed, total receipts are EUR 1.25 million and payments EUR 1.6 million, producing a net short EUR 0.35 million. Even the sign can change.

When credible capture bounds are available, they yield transparent partial identification. If $s^+\in[\underline s^+,\overline s^+]$ and $s^-\in[\underline s^-,\overline s^-]$, then
\begin{equation}
\frac{Y^+}{\overline s^+}-\frac{Y^-}{\underline s^-}
\ \leq F^+-F^-\leq\
\frac{Y^+}{\underline s^+}-\frac{Y^-}{\overline s^-}.
\label{eq:capturebounds}
\end{equation}
The lower bound combines the smallest feasible receipts with the largest feasible payments; the upper bound reverses them. Correlated or jointly restricted shares may tighten this interval. Without such restrictions the wide interval is information about uncertainty, not a defect to conceal through a more flexible model.

Additional evidence might include client-provided treasury forecasts, authorized consolidated statements, verified receivables and payables, or known multi-bank account data. Financial statements can constrain annual totals but usually do not reveal monthly currency exposure. Web information can classify business activities or identify a dated event; it rarely identifies bank shares. Until reliable additional observations exist, sales explanations should say ``activity visible to this bank'' and invite verification of total exposure.

\section{Missingness, zeros, and dates}
A zero means no qualifying transaction was observed in a complete month. A missing record means that the bank does not know whether activity occurred. Account opening, migration, closure, and failed ingestion create different missingness mechanisms. A model that replaces each missing value with zero teaches itself that clients become inactive when data disappear.

Let $O_{ict}^{\pm}$ indicate that a flow is observed. The likelihood conditions on the recorded observation pattern when missingness is ignorable given available information. If closure depends on unobserved activity or a deteriorating relationship, this assumption may fail. A closure or capture model can then be added, but separate identification will depend on extra signals. The first product should expose observation coverage and use a fallback for incomplete clients.

Every feature needs both an economic date and an availability date. GDP is typically released at quarterly frequency and revised later; carrying a released quarterly measure forward to monthly decision dates is defensible, while interpolating future unreleased GDP into past months leaks information. The same rule applies to revised industry statistics, financial statements, website content, and future sales outcomes. A present-day website cannot be treated as a historical signal unless its historical availability is established.

Address and industry data help assign peers and currency-market exposure. Names are useful for entity resolution but have little economic meaning as arbitrary high-cardinality identifiers. Industry classifications, country, business size, relationship tenure, products already held, and verified export/import activity are more interpretable predictors. New or revised attributes must be dated, and uncertainty in entity matching should remain visible.

\section{The data contract for the separate modules}
The forecasting module requires native-currency receipt and payment series, observation masks, client attributes, macroeconomic vintages, and a forecast origin. Its output is a set of weighted joint paths with currency and date labels, together with a declared coverage scope: bank-observed or total-wallet under stated capture assumptions.

The FX module additionally requires existing cash balances, known contractual commitments, hedge positions, netting rules, settlement dates, and product terms. The choice module requires the actual alternatives offered or available, timestamps, quoted forward points and fees, eligibility, acceptance or decline, and the outcome-observation window. The outreach module needs assignment or contact records, the information available before assignment, salesperson capacity, and realized net contribution.

These contracts allow useful partial delivery. A sound flow forecast need not wait for full preference data. Conversely, the existence of transaction histories does not imply that a choice model is identified. Logging the missing offers and outcomes can begin while the forecasting system is developed.

\chapter{Learning dynamics from many short histories}
\label{ch:dynamics}
\section{Pooling is the starting point}
The forecasting literature offers a useful change in perspective: the bank has a panel-learning problem rather than one million unrelated time-series problems. \citet{montero2021global} distinguish local models, estimated separately for each series, from global models that share a forecasting rule. Their Proposition~1 establishes an expressiveness equivalence for functions on complete finite histories under its stated construction. That mathematical result does not say that an arbitrary small global network reproduces every local rule, nor that empirical generalization follows without restrictions. Finite-memory inputs that give the same input to two clients with different desired outputs illustrate why representation matters.

A hierarchical regression makes the useful statistical mechanism explicit. For one client, let $\bm y_i=X_i\beta_i+\varepsilon_i$ with $\varepsilon_i\sim\N(0,\sigma^2 I)$. Rows of $X_i$ can contain lagged flows and known calendar or macro inputs. Assume
\begin{equation}
\beta_i\mid g(i)\sim\N(\beta_{g(i)},\Lambda^{-1}),
\end{equation}
where $g(i)$ is an industry or peer group and $\Lambda$ is prior precision. Combining the Gaussian likelihood and prior yields the negative log posterior, up to constants,
\begin{equation}
\frac{1}{2\sigma^2}(\bm y_i-X_i\beta_i)'(\bm y_i-X_i\beta_i)
+\frac12(\beta_i-\beta_g)'\Lambda(\beta_i-\beta_g).
\end{equation}
Differentiating and setting the derivative to zero gives
\begin{align}
(X_i'X_i/\sigma^2+\Lambda)\widehat\beta_i
 &=X_i'\bm y_i/\sigma^2+\Lambda\beta_g,\\
\widehat\beta_i
 &=(X_i'X_i/\sigma^2+\Lambda)^{-1}
 (X_i'\bm y_i/\sigma^2+\Lambda\beta_g),
\label{eq:shrinkage}\\
\V(\beta_i\mid\bm y_i)&=(X_i'X_i/\sigma^2+\Lambda)^{-1}.
\end{align}
When a client has little informative history, the peer estimate carries more weight. With a well-conditioned, highly informative client history, its own data dominate. In the scalar case with information $s=X_i'X_i/\sigma^2$ and prior precision $\lambda$, the posterior mean is $s/(s+\lambda)$ times the local estimate plus $\lambda/(s+\lambda)$ times the peer mean. This is the precise sense in which pooling can help a three-year series. Equation~\eqref{eq:shrinkage} is our Gaussian derivation, motivated by the global-versus-local question; it is not a theorem attributed to the global-forecasting paper.

Group definitions and shrinkage strengths should be estimated or validated, not assumed harmless. Pooling exporters and importers without currency and business-type information can average away the very differences needed for FX advice. A practical hierarchy shares most parameters while permitting industry and client deviations only where observations support them.

\section{A latent state gives the forecast a financial interpretation}
\label{sec:kalman}
Cash receipts vary because persistent business conditions change and because individual payments arrive irregularly. A state-space model separates those components. Let $x_t\in\mathbb R^d$ be a latent state such as current activity and slowly changing seasonal conditions. Temporarily suppress client indices and assume a linear Gaussian system:
\begin{align}
x_t&=A_t x_{t-1}+b_t+\eta_t,&\eta_t&\sim\N(0,Q_t),\\
y_t&=H_t x_t+d_t+\epsilon_t,&\epsilon_t&\sim\N(0,R_t).
\label{eq:ssm}
\end{align}
All innovations are mutually independent over time and between equations, and $x_0\sim\N(m_0,P_0)$. The offsets $b_t,d_t$ can contain known macroeconomic and calendar inputs. The vector $y_t$ can represent transformed inflows and outflows together. A positive off-diagonal covariance can express their shared business cycle.

\citet{kalman1960} derives recursive linear estimation using orthogonality and covariance propagation, particularly the main filtering theorem and covariance recursion. Under the Gaussian assumptions above, the linear estimator is the full conditional mean and the conditional distribution remains Gaussian. Here is the derivation in the notation needed for the bank.

Suppose $x_{t-1}\mid y_{1:t-1}\sim\N(m_{t-1},P_{t-1})$. Applying the affine transition gives the prior at the next date:
\begin{equation}
a_t=A_t m_{t-1}+b_t,\qquad
P_t^-=A_tP_{t-1}A_t'+Q_t.
\end{equation}
The forecast error is $v_t=y_t-H_ta_t-d_t$, with covariance $S_t=H_tP_t^-H_t'+R_t$. The joint conditional distribution of $x_t$ and $v_t$ has mean $(a_t,0)$ and block covariance
\begin{equation}
\begin{pmatrix}
P_t^- & P_t^-H_t'\\
H_tP_t^- & S_t
\end{pmatrix}.
\end{equation}
To obtain the conditional mean, subtract $K_tv_t$ from $x_t-a_t$, where $K_t=P_t^-H_t'S_t^{-1}$. Its covariance with $v_t$ is $P_t^-H_t'-K_tS_t=0$. Joint Gaussianity turns zero covariance into independence. Therefore
\begin{align}
m_t&=a_t+K_tv_t,\\
P_t&=P_t^- -P_t^-H_t'S_t^{-1}H_tP_t^-.
\label{eq:filter}
\end{align}
The gain $K_t$ gives less weight to a noisy observation and more weight to one that sharply resolves prior uncertainty. In a scalar example with prior variance $4$ and observation variance $1$, $K=4/5$: an observed log-flow surprise of $0.10$ raises the estimated state by $0.08$.

The forecast-error distribution supplies the exact Gaussian likelihood:
\begin{equation}
\ell(\theta)=-\frac12\sum_t
\left[k_t\log(2\pi)+\log|S_t|+v_t'S_t^{-1}v_t\right],
\label{eq:kalmanlik}
\end{equation}
where $k_t$ is the number of observed components. Missing components are removed by a selection matrix before computing $v_t,S_t$; missing all components means skipping the update. Linear solves and factorizations should replace explicit inverses in software. The equivalent Joseph covariance update,
\begin{equation}
P_t=(I-K_tH_t)P_t^-(I-K_tH_t)'+K_tR_tK_t',
\end{equation}
helps preserve covariance symmetry and positive semidefiniteness in finite precision. The formula changes the numerical evaluation, not the statistical target.

For fixed future coefficients, repeated propagation gives
\begin{equation}
\E[x_{t+h}\mid\Hist_t]=A^h m_t+\sum_{r=0}^{h-1}A^r b_{t+h-r},
\end{equation}
and, with constant $A,Q$,
\begin{equation}
\V(x_{t+h}\mid\Hist_t)=A^hP_t(A^h)'+
\sum_{r=0}^{h-1}A^rQ(A^r)'.
\label{eq:multistep}
\end{equation}
Sampling one state path recursively preserves cross-horizon dependence. Drawing each month independently from its marginal distribution loses that dependence and can misstate cumulative exposure risk.

\section{Smoothing and EM: how the parameters can be learned}
\label{sec:em}
Filtering uses observations only up to each date. Parameter estimation can use the whole training window to infer past states. \citet{shumway1982} combine smoothing with expectation--maximization (EM); their equations (3)--(4) give the complete-data likelihood and its conditional expectation, and their Appendix gives filtering, smoothing, and lag-covariance recursions.

For a time-invariant transition $x_t=Ax_{t-1}+\eta_t$, the filtered joint distribution of $(x_t,x_{t+1})$ implies
\begin{equation}
\E[x_t\mid x_{t+1},y_{1:t}]
 =m_t+J_t(x_{t+1}-a_{t+1}),\qquad
J_t=P_tA'(P_{t+1}^-)^{-1}.
\end{equation}
The conditional Markov property means that future observations add information about $x_t$ through $x_{t+1}$. Taking expectations and variances conditional on all training observations gives the backward recursion
\begin{align}
\widetilde m_t&=m_t+J_t(\widetilde m_{t+1}-a_{t+1}),\\
\widetilde P_t&=P_t+J_t(\widetilde P_{t+1}-P_{t+1}^-)J_t'.
\label{eq:smoothing}
\end{align}
Start with $\widetilde m_T=m_T$ and $\widetilde P_T=P_T$. The same conditional decomposition gives
\begin{equation}
\Cov(x_t,x_{t-1}\mid y_{1:T})=\widetilde P_tJ_{t-1}'.
\label{eq:lagcov}
\end{equation}
Indeed, $x_{t-1}$ equals an affine function with slope $J_{t-1}$ of $x_t$, plus a residual independent of $x_t$ and subsequent observations. Multiplying that relation by the centered $x_t$ yields~\eqref{eq:lagcov}. These moments, not products of smoothed means alone, are needed for unbiased complete-data sufficient statistics.

Define sums over the $T$ transitions:
\begin{align}
S_{00}&=\sum_t(\widetilde P_{t-1}+\widetilde m_{t-1}\widetilde m_{t-1}'),\\
S_{10}&=\sum_t(\widetilde P_tJ_{t-1}'+\widetilde m_t\widetilde m_{t-1}'),\\
S_{11}&=\sum_t(\widetilde P_t+\widetilde m_t\widetilde m_t').
\end{align}
The part of the expected complete-data log likelihood involving $A,Q$ is
\begin{equation}
-\frac{T}{2}\log|Q|-\frac12\tr\left[
Q^{-1}(S_{11}-AS_{10}'-S_{10}A'+AS_{00}A')\right].
\end{equation}
Differentiating with respect to $A$ gives $AS_{00}=S_{10}$. Substitution and differentiation with respect to $Q$ yield
\begin{equation}
A^{\rm new}=S_{10}S_{00}^{-1},\qquad
Q^{\rm new}=\frac1T(S_{11}-S_{10}S_{00}^{-1}S_{10}').
\label{eq:emupdate}
\end{equation}
For a known observation matrix $H$ and complete observations,
\begin{equation}
R^{\rm new}=\frac1T\sum_t
\left[(y_t-H\widetilde m_t)(y_t-H\widetilde m_t)'
      +H\widetilde P_tH'\right].
\end{equation}
Offsets are subtracted first. Exogenous regressors can be appended to the transition regression with their deterministic cross-moments. Hierarchical penalties and structural restrictions generally replace these unrestricted closed-form updates with constrained optimization. For a diagonal $R$ and missing observations, each diagonal update uses its own observed count; a general unstructured $R$ with changing missingness requires the appropriate observed-data objective, not an indiscriminate division by $T$.

Why does exact EM improve likelihood? For any density $q(x)$,
\begin{equation}
\log p_\theta(y)
=\underbrace{\E_q[\log p_\theta(y,x)-\log q(x)]}_{\mathcal L(q,\theta)}
+\KL\{q(x)\Vert p_\theta(x\mid y)\}.
\label{eq:elbo}
\end{equation}
This follows by expanding the KL divergence and using
$\log p_\theta(x\mid y)=\log p_\theta(y,x)-\log p_\theta(y)$.
The E-step sets $q$ to the old posterior, making the KL term zero at the old parameter. Maximizing $\mathcal L$ over $\theta$ cannot lower the observed likelihood because the new KL term is nonnegative. Approximate smoothing removes this exact argument, and even exact EM can converge to a local stationary solution rather than a global optimum. Both distinctions matter for the non-Gaussian extension below.

At bank scale, a low-dimensional filter costs roughly $O(Td^3)$ per client in a dense implementation, with observation-dimension costs in addition. The advantage is separability across clients conditional on shared factors and parameters. A million separate high-dimensional covariance models would waste data and computation. Shared low-rank covariance structures and small client states are the appropriate starting hypotheses.

\section{Macroeconomic factors compress predictors but do not identify causes}
\citet{stock2002} study forecasting with principal components from many predictors. Their Sections~2--3 set out the approximate factor representation and forecasting results under assumptions that let both the number of predictors and the time dimension grow. That asymptotic setting does not describe 36 dates merely because the bank has many clients.

Let $X$ be a centered, standardized $T\times K$ matrix of macroeconomic predictors, with rows denoting dates. An approximate rank-$r$ factor representation solves
\begin{equation}
\min_{F,\Lambda}\|X-F\Lambda'\|_F^2
\quad\text{subject to}\quad F'F/T=I_r.
\end{equation}
For fixed $F$, differentiating gives $\widehat\Lambda=X'F/T$. Substitution leaves a constant minus $\tr(F'XX'F)/T$. The maximizing $F$ consists of $\sqrt T$ times the top $r$ orthonormal eigenvectors of $XX'$. Forecasting then regresses future outcomes on these estimated factors, their lags, and suitable outcome lags.

The compression is useful when many macro indicators measure overlapping conditions. It also has two limitations. First, $(F,\Lambda)$ and $(FR,\Lambda R)$ give the same fitted $X$ for an orthogonal $R$; individual factor names require additional interpretation. Second, predictive factors can incorporate movements also associated with GDP. If a free factor absorbs the GDP path, a separately named GDP loading need not be identified. A proposal requiring an interpretable GDP response should put that variable explicitly in the model and restrict residual factors, for example by a declared projection or structural exclusion. Orthogonalization defines a predictive decomposition; it does not prove causal exogeneity.

Longer external macro histories can help estimate factor dynamics or prior scenario distributions. They cannot reveal a particular client's response before the bank observed that client. Vintage-consistent estimation and sensitivity to the number of factors remain necessary. With a short bank history, a small set of economically chosen exogenous inputs is the primary specification; principal components are a challenger when their incremental forecast value can be demonstrated.

\section{Two neural alternatives and what they contribute}
\subsection{DeepAR: share a distributional forecasting rule}
\citet{salinas2020deepar} construct a global autoregressive recurrent network. Their model and likelihood sections define a hidden state updated from the previous target and current covariates,
\begin{equation}
h_{it}=f_\theta(h_{i,t-1},y_{i,t-1},x_{it}),\qquad
\vartheta_{it}=g_\theta(h_{it}),
\end{equation}
and factor the predictive likelihood as
\begin{equation}
p_\theta(y_{i,1:T}\mid x_i)=\prod_t p(y_{it}\mid\vartheta_{it}).
\label{eq:deepar}
\end{equation}
Training maximizes the sum of log factors across clients and time. During training the lagged targets are observed; future generation replaces them with recursively sampled targets. This distinction is essential because feeding future actual observations at evaluation would overstate forecast quality.

For real-valued observations, a Gaussian output can use a positive scale $\sigma=\log(1+e^a)$. For nonnegative counts the paper also uses a negative binomial distribution. With mean $\mu>0$ and dispersion $\alpha>0$, set $r=1/\alpha$ and $p=r/(r+\mu)$, giving
\begin{equation}
\Pr(Y=y)=\frac{\Gamma(y+r)}{\Gamma(r)\Gamma(y+1)}
p^r(1-p)^y,\quad
\E[Y]=\mu,\quad\V(Y)=\mu+\alpha\mu^2.
\end{equation}
The moments follow from the probability-generating function
$G(s)=[p/(1-(1-p)s)]^r$: $G'(1)=r(1-p)/p$ and
$G''(1)+G'(1)-G'(1)^2=r(1-p)/p^2$.
These count likelihoods are appropriate for numbers of transactions, not directly for dollar or euro amounts. Continuous positive amounts need a continuous distribution, with a separate point mass at zero if necessary.

For this project, DeepAR is a meaningful flexible comparator and a source of shared distributional learning. Its basic per-series likelihood does not itself provide correlated inflow/outflow paths, a common macroeconomic shock across firms, or a structurally interpretable GDP response. Those would require an explicit multivariate or shared-factor extension. The published benchmarks establish performance in their data settings, not in this bank's partial-wallet setting. We therefore retain a structured model as the proposed first release and make neural complexity earn its place in held-out decision-relevant tests.

\subsection{Deep state space: share parameters while retaining filtering}
\citet{rangapuram2018} combine a recurrent parameter generator with a linear Gaussian state-space model. In their notation the latent state $\ell_t$ evolves linearly and observations are linear in the preceding state; an index shift maps their equations (1)--(4) to~\eqref{eq:ssm}. The shared recurrent network maps known covariates to time-varying transition, noise, and observation parameters. Crucially, in their construction the observations update the latent state through filtering; they are not inputs to the parameter-generating RNN. Conditional on the covariates and network parameters, the likelihood is the analytic Kalman likelihood~\eqref{eq:kalmanlik}.

Writing $\theta_t=\Psi_\phi(x_{1:t})$ for these generated parameters, estimation maximizes
\begin{equation}
\sum_i\log p(y_{i,1:T}\mid\theta_{i,1:T})
\end{equation}
by differentiating through the filter. This creates two forms of adaptation: shared neural weights learn how covariates should change dynamics, and each client's filter updates its own state from observed transactions. The source assumes conditional independence across series in this objective. It therefore motivates shared estimation, while the bank's desired dependence across currencies and firms is an additional modeling responsibility.

The official long version's Appendix~A.2 describes parameter constraints, and Appendix~A.4 sketches a variational extension to non-Gaussian observations. The latter does not make arbitrary zero-inflated observations exactly Gaussian-filterable. For this proposal, positivity of variances, stable or deliberately integrated states, and monotonic scenario effects must be designed explicitly. A recurrent parameter generator is optional after the transparent hierarchy and hurdle likelihood have been validated. Neither a neural transition nor an exogenous input label, by itself, gives a causal counterfactual interpretation.

\chapter{The proposed joint inflow--outflow model}
\label{ch:joint}
\section{A shared dynamic hurdle model}
Cash flows have two features that the Gaussian teaching model does not capture directly: many currency-months are exactly zero, and positive amounts are skewed. We propose a dynamic hurdle model with a shared latent state. The hurdle determines whether a positive flow occurs; a continuous distribution determines its size. This is a proposed synthesis of the pooling and state-space methods already derived, rather than a claim that one cited paper supplies the entire banking model.

Use $r=(c,s)$ to index currency $c$ and direction $s\in\{+,-\}$. Let $Y_{irt}$ denote the bank-observed flow in that coordinate and $D_{irt}=\ind\{Y_{irt}>0\}$. For an observed coordinate, define
\begin{align}
D_{irt}\mid x_{it},f_t&\sim\operatorname{Bernoulli}(p_{irt}),
&p_{irt}&=\sigmoid(\eta^D_{irt}),\\
\log(Y_{irt}/a_r)\mid D_{irt}=1,x_{it},f_t
&=\eta^A_{irt}+\varepsilon_{irt},
&\bm\varepsilon_{it}&\sim\N(0,R_{g(i)}).
\label{eq:hurdle}
\end{align}
Here $a_r$ is a fixed positive unit scale, such as one thousand units of the native currency, chosen using training data or a declared unit convention. It avoids taking the logarithm of a dimensional quantity. The logistic function is $\sigmoid(v)=1/(1+e^{-v})$. Conditional on the states, hurdle indicators are initially independent across coordinates and independent of the Gaussian magnitude errors. Correlated latent states still create marginal dependence in occurrence. If simultaneous zero/nonzero patterns remain poorly described, correlated occurrence shocks are a targeted extension.

The local state $x_{it}$ and common state $f_t$ obey
\begin{align}
x_{it}&=A_{g(i)}x_{i,t-1}+\eta_{it},&\eta_{it}&\sim\N(0,Q_{g(i)}),\\
f_t&=A_f f_{t-1}+\xi_t,&\xi_t&\sim\N(0,Q_f).
\end{align}
Local and common innovations are independent in the initial specification. Common states may be global or sector-specific, with a small fixed dimension selected by validation. A client state can encode a persistent activity level and deviations from shared seasonality. For each of occurrence and magnitude, the linear predictor is
\begin{equation}
\eta^k_{irt}=\alpha^k_{g(i),r}+z_i'\gamma^k_r+u^k_{ir}
+b^k_{r}(t)+{h^k_r}'x_{it}+{\lambda^k_{ir}}'f_t
+B^k_{ir}(m_{1:t}),\qquad k\in\{D,A\}.
\label{eq:predictor}
\end{equation}
The term $b_r^k(t)$ uses a finite set of known calendar predictors, such as month-of-year and holiday indicators; $u^k_{ir}$ is a shrunk client deviation; and $B^k_{ir}$ is the explicit macroeconomic response, constructed in Chapter~\ref{ch:cf}. Unrestricted date effects would absorb a common GDP path, so they are excluded when estimating that response. Residual common factors still require the separation restrictions discussed in Chapter~\ref{ch:dynamics}; shrinkage alone does not identify a causal GDP coefficient. Missing or sparse currency coordinates borrow information through the hierarchy rather than receiving unrestricted coefficients. An arbitrary client identifier may index a random effect but is not treated as an economically meaningful covariate.

Separation between intercepts and persistent states requires a normalization. One defensible specification uses zero-mean stationary local deviations and a separate client intercept; another uses a local level with a fixed initial-mean hierarchy and no redundant free intercept. We recommend the former baseline. Stationary transitions require eigenvalues inside the unit circle; an intentional random-walk level is a different hypothesis with growing forecast uncertainty. Seasonal coefficients can satisfy a sum-to-zero constraint. Factor scale and rotation need normalization, for example a constrained loading block. These restrictions remove representational ambiguity, not economic uncertainty.

\section{The likelihood and the dependence it preserves}
Let $\mathcal O_{it}$ be the observed coordinates and $\mathcal P_{it}$ the subset with positive amounts. Write $\bm v_{it}$ for the observed log magnitudes in $\mathcal P_{it}$, and $R_{\mathcal P\mathcal P}$ for the corresponding covariance submatrix. Conditional on the latent states, the likelihood contribution is
\begin{align}
p(\bm Y_{it}\mid x_{it},f_t)
={}&\prod_{r\in\mathcal O_{it}}p_{irt}^{D_{irt}}(1-p_{irt})^{1-D_{irt}}\nonumber\\
&\times\phi_{|\mathcal P_{it}|}
 (\bm v_{it};\bm\eta^A_{\mathcal P_{it}},R_{\mathcal P\mathcal P})
 \prod_{r\in\mathcal P_{it}}Y_{irt}^{-1}.
\label{eq:hurdlelik}
\end{align}
The Jacobian $1/Y$ converts a Gaussian log-amount density to an amount density. If there are no positive amounts, the Gaussian factor is one. Zero coordinates do not contribute a Gaussian observation of zero log amount. Missing coordinates contribute neither a Bernoulli factor nor an amount density. The covariance submatrix arises by marginalizing unobserved Gaussian log magnitudes; it is not a conditional covariance given invented zero magnitudes.

For one coordinate, conditional on the states,
\begin{align}
\E[Y_{irt}\mid x,f]&=a_r p_{irt}
 \exp(\eta^A_{irt}+R_{rr}/2),\label{eq:hurdlemean}\\
\E[Y_{irt}^2\mid x,f]&=a_r^2p_{irt}
 \exp(2\eta^A_{irt}+2R_{rr}),\\
\V(Y_{irt}\mid x,f)&=\E[Y_{irt}^2\mid x,f]-\E[Y_{irt}\mid x,f]^2.
\end{align}
These equations follow from $D^2=D$, conditional independence, and the normal moment-generating function. For distinct coordinates $r,s$ under the specified conditional occurrence independence,
\begin{equation}
\E[Y_rY_s\mid x,f]
=a_ra_s p_rp_s\exp\left(\eta_r^A+\eta_s^A+
\frac{R_{rr}+R_{ss}}{2}+R_{rs}\right).
\end{equation}
Thus correlated magnitude noise and shared latent conditions both affect joint exposure. Integrating over states requires $\E_{x,f}[p(x,f)e^{\eta^A(x,f)+R/2}]$; replacing it by the product of separate expectations is generally wrong because occurrence and magnitude depend on the same states.

For same-currency net flow $N=Y^+-Y^-$,
\begin{equation}
\V(N)=\V(Y^+)+\V(Y^-)-2\Cov(Y^+,Y^-).
\label{eq:netvariance}
\end{equation}
If both standard deviations are EUR 100,000 and correlation is $0.9$, the net standard deviation is about EUR 44,721. Assuming independence gives EUR 141,421. The difference changes hedge sizing and the value of risk reduction. A model that fits marginal means well can still be unsuitable for hedging if it misses this covariance.

Across clients, conditional independence given $f$ remains computationally useful. Unconditionally,
\begin{equation}
\Cov(Y_i,Y_j)=
\E[\Cov(Y_i,Y_j\mid f)]+
\Cov(\E[Y_i\mid f],\E[Y_j\mid f]).
\end{equation}
The first term is zero under the baseline independence assumptions for different clients, but the second need not be. Shared draws of future $f$ create coherent sector or portfolio shocks. Drawing a separate common factor for every client would remove the intended dependence.

\section{Inference for zeros and positive amounts}
The linear Gaussian model admits exact filtering. The hurdle model does not. A practical first implementation can use a Gaussian approximation to the latent posterior, with exact Gaussian calculations retained as a reference case.

For one client conditional on common factors, stack states into $\bm x=(x_0',\ldots,x_T')'$. Its Gaussian Markov prior has a block-tridiagonal precision matrix $K$. For fixed covariance and loadings, the log posterior is a sum of the Gaussian prior, concave Bernoulli log likelihoods, and quadratic Gaussian log-magnitude terms:
\begin{align}
L(\bm x)={}&-\tfrac12(\bm x-\bm\mu)'K(\bm x-\bm\mu)\nonumber\\
&+\sum_{r,t\in\mathcal O}\{D_{rt}\eta^D_{rt}-\log(1+e^{\eta^D_{rt}})\}\nonumber\\
&-\tfrac12\sum_t(\bm v_t-\bm\eta_t^A)'R_{\mathcal P\mathcal P}^{-1}
                  (\bm v_t-\bm\eta_t^A)+\text{constant}.
\label{eq:laplaceobj}
\end{align}
With affine predictors, a Bernoulli term has gradient $(D-p)h$ and negative Hessian $p(1-p)hh'$. The Gaussian magnitude term contributes $H'R^{-1}H$. Hence the negative Hessian is the positive-definite prior precision plus positive-semidefinite local terms. With a proper prior this conditional posterior is log-concave, and a safeguarded Newton method can find its unique mode. If $\widehat x$ is that mode and $J=-\nabla^2L(\widehat x)$, the Laplace approximation is
\begin{equation}
q(\bm x)=\N(\widehat x,J^{-1}),\qquad
\log p(y)\approx L(\widehat x)+\tfrac{d(T+1)}2\log(2\pi)-\tfrac12\log|J|,
\end{equation}
where $L$ must include the likelihood and prior normalizing constants for the marginal-likelihood formula. Keeping the band structure avoids a dense inverse of a $d(T+1)$ matrix. Nonlinear predictors or simultaneous parameter optimization can destroy this simple concavity; the statement is conditional on fixed linear-Gaussian structural parameters.

An alternative is structured variational inference. For a Gaussian Markov family $q_\phi(x)$, maximize the lower bound $\mathcal L(q_\phi,\theta)$ in~\eqref{eq:elbo}. Write $x=\mu_\phi+L_\phi\epsilon$ with $\epsilon\sim\N(0,I)$ to differentiate Monte Carlo estimates of the expected log likelihood. The Gaussian prior and entropy terms can be evaluated analytically. A diagonal variational family is computationally convenient but discards temporal posterior covariance needed for cumulative exposure; a banded precision or state-space variational family is preferable. The variational sketch in \citet[Appendix A.4]{rangapuram2018} motivates this route, while the likelihood in~\eqref{eq:hurdlelik} and the proposed hierarchy are specific to the current problem.

Approximate EM based on Laplace moments is a feasible engineering option, but it inherits neither exact EM monotonicity nor guaranteed interval calibration. Before scale-up, compare small problems against numerical integration or a carefully checked posterior sampler, and compare the continuous Gaussian special case against the exact filter. These are validity checks of approximation error. Predictive backtests then answer the separate question of forecasting usefulness.

\section{A feasible estimation and simulation sequence}
Begin with a transparent Gaussian transformed-flow model and shared seasonal effects as an implementation reference, then fit the hurdle likelihood on the same chronological splits. Estimate shared and group coefficients from many clients, using client-specific states and shrunk deviations. Update common factors using their joint contribution across clients, rather than plugging in a noiseless cross-sectional mean. A practical variational factorization may alternate a common-factor posterior with client-conditional posteriors; any lost dependence must be measured on a smaller joint reference problem. With only 36 dates the common-factor path is short, although each date receives many observations.

For every forecast origin, simulate complete paths in this order: draw shared parameters or a parameter-uncertainty approximation; draw common future factors once for the scenario; draw each client's current state from its posterior; propagate its states; draw occurrence indicators and correlated positive log amounts; and convert to native-currency inflows and outflows. Market FX paths must be linked to the same macro scenario and, where economically material, to common innovations. Independent FX simulation is an explicit simplifying assumption, not the default meaning of a joint forecast.

Parameter uncertainty matters particularly for macro slopes and sparse currencies. Plug-in predictions capture future process noise while omitting uncertainty about estimated coefficients. A hierarchical posterior, suitably clustered parameter bootstrap, or a declared sensitivity envelope can address part of this gap. No method removes the limited evidence about macroeconomic regimes absent from the training period.

The output to the utility module is a reproducible collection
\begin{equation}
\{w_s,\;Y^{+,(s)}_{ic,t+h},\;Y^{-,(s)}_{ic,t+h},\;S^{(s)}_{c,t+h}\}_{s=1}^{M},
\qquad w_s\geq0,\quad\sum_s w_s=1,
\end{equation}
with units, dates, scenario name, as-of information date, capture assumptions, and model version. Weights describe the probability measure used. Physical probabilities are appropriate for expected client outcomes; market pricing of derivative terms uses a separate pricing measure or observed quotes. The next chapters explain why those distinctions must survive the interface.

\chapter{Macroeconomic scenarios and genuine counterfactuals}
\label{ch:cf}
\section{Three distinct questions}
The time-series model should answer how future receipts and payments change when macroeconomic conditions change. Giving GDP a coefficient is a necessary piece of that design, but the interpretation of the change needs care. \citet[Sections 3.2--3.4]{pearl2009} distinguishes observing a variable, intervening on its generating mechanism, and reasoning about a counterfactual for a particular observed unit. These are different probability statements:
\begin{align}
&p(Y_{t+1:t+H}\mid\Hist_t,M_{t+1:t+H}=m)
&&\text{conditional forecast},\label{eq:conditional}\\
&p(Y_{t+1:t+H}\mid\Hist_t,\doop(M_{t+1:t+H}=m))
&&\text{intervention distribution},\label{eq:intervention}\\
&p(Y^{m'}_{t+1:t+H}\mid\Hist_t,\text{observed factual evidence})
&&\text{unit counterfactual}.\label{eq:unitcf}
\end{align}
Here $Y^{m'}$ is the outcome in an alternative structural system, using information about the latent conditions of the factual entity. A bank may use an externally supplied GDP scenario without claiming that it or any single policymaker can directly set GDP. If the scenario is implemented by inserting $m$ into a forecasting equation while keeping residual innovations at their unconditional distribution, that operation equals~\eqref{eq:conditional} only under the corresponding exogeneity assumptions. It is otherwise a model-based forcing scenario.

For an ordinary forecast in which future macroeconomic conditions are unknown, the corresponding predictive distribution integrates over an as-of-date macro forecast:
\begin{equation}
p(Y_{t+1:t+H}\mid\Hist_t)=\int p(Y_{t+1:t+H}\mid\Hist_t,m)\,p(m\mid\Hist_t)\,dm.
\end{equation}
This is the law of total probability for a coherent joint model. Fixing one macro path omits uncertainty across macro paths. A stress scenario may instead be deliberately assigned no probability; it remains a conditional financial calculation rather than an unconditional risk forecast.

The distinction can be proved with a small linear example. Suppose
\begin{equation}
M=U+\nu,\qquad Y=\beta M+\lambda U+\varepsilon,
\end{equation}
where independent centered Gaussian variables have variances $\sigma_U^2,\sigma_\nu^2,\sigma_\varepsilon^2$. Because
$\Cov(U,M)=\sigma_U^2$ and $\V(M)=\sigma_U^2+\sigma_\nu^2$,
\begin{equation}
\E[Y\mid M=m]=
\left(\beta+\lambda\frac{\sigma_U^2}{\sigma_U^2+\sigma_\nu^2}\right)m.
\label{eq:confounding}
\end{equation}
An intervention replaces the first equation by $M=m$ and leaves $U$ distributed as before. Consequently,
\begin{equation}
\E[Y\mid\doop(M=m)]=\beta m.
\end{equation}
The extra slope in~\eqref{eq:confounding} is not a GDP effect; it comes from the common cause. More clients measured at the same dates do not resolve this confounding.

\section{Structural assumptions needed for macroeconomic interventions}
A structural cash-flow system can be written
\begin{align}
M_t&=g_M(M_{<t},U^M_t),\\
X_{it}&=g_X(X_{i,t-1},M_{\leq t},Z_i,U^X_{it}),\\
F^\pm_{ict}&=g^\pm_F(X_{it},M_{\leq t},Z_i,U^F_{ict}),\\
Y^\pm_{ict}&=g^\pm_O(F^\pm_{ict},S^\pm_{ict},U^O_{ict}).
\end{align}
The variables $S^\pm$ here denote bank capture, not exchange rates. Exogenous disturbances may share dependence unless assumptions exclude it. An intervention on $M$ replaces its structural equation and propagates through the others. The response has causal meaning only if the retained equations and disturbance distributions are invariant to that intervention and their relevant parameters are identified or explicitly assumed.

For a single exposure $M$, if observed $Z$ blocks all back-door confounding paths and contains no inappropriate descendants, the adjustment formula is
\begin{equation}
p(y\mid\doop(m))=\int p(y\mid m,z)p(z)\,dz.
\label{eq:backdoor}
\end{equation}
The derivation uses conditional exchangeability $Y(m)\perp M\mid Z$, consistency $Y=Y(M)$, and positivity. First integrate $p(Y(m)=y\mid Z=z)$ over $p(z)$; exchangeability permits conditioning additionally on $M=m$; consistency then replaces $Y(m)$ by the observed $Y$. Longitudinal interventions require the corresponding sequential assumptions and time ordering, not an undated regression adjustment. Pearl's back-door discussion gives the graphical criteria behind the statistical assumptions.

In this banking application, candidate confounders include sector shocks, policy responses, commodity prices, and exchange-rate conditions. GDP is an aggregate outcome of many mechanisms, so a ``one percentage point GDP increase'' can correspond to several economically different interventions. A demand stimulus and a productivity improvement can raise GDP while affecting margins, payments, and FX rates differently. The first release should therefore label scenarios by the variables set, accompanying assumptions, and whether the response is predictive, imposed, or causally identified. Identifying a particular policy shock would be a separate econometric project.

\section{How to impose the requested positive GDP response}
The user requests a scenario in which higher GDP increases every entity's mean cash inflow. The model can satisfy this transparently. Let $g_t$ be a declared real-GDP growth measure expressed in percentage points, centered on the training mean. For inflow magnitude, use a finite distributed lag
\begin{equation}
B^A_{ic+}(m_{1:t})
=\beta^A_{ic}\sum_{\ell=0}^{L}w^A_{c\ell}g_{t-\ell}
+\widetilde B^A_{ic}(m_{1:t}^{\rm other}),
\label{eq:gdpkernel}
\end{equation}
with
\begin{equation}
\beta^A_{ic}=\softplus(b^A_{g(i),c}+z_i'\gamma^A_c)>0,
\quad w^A_{c\ell}\geq0,\quad\sum_{\ell=0}^{L}w^A_{c\ell}=1.
\end{equation}
The lag normalization separates total response size $\beta$ from timing. Weights may be shared across industries or currencies; a softmax is one possible positive parameterization. An analogous occurrence term uses $\beta^D_{ic}\geq0$ and nonnegative weights, or sets that channel to zero if positivity is imposed only on amount. Exact zero responses can be permitted through a nonnegative constrained parameter rather than strict softplus.

For fixed latent states, fixed residual variance, and a GDP input $g_{t-\ell}$, differentiating the conditional mean in~\eqref{eq:hurdlemean} gives
\begin{align}
\frac{\partial\log\E[Y^+_{ict}\mid x,f,m]}{\partial g_{t-\ell}}
&=\frac{\partial\log p_{ict}}{\partial\eta^D_{ict}}
      \beta^D_{ic}w^D_{c\ell}+\beta^A_{ic}w^A_{c\ell}\nonumber\\
&=(1-p_{ict})\beta^D_{ic}w^D_{c\ell}
       +\beta^A_{ic}w^A_{c\ell}\geq0.
\label{eq:monotone}
\end{align}
The identity $\partial\log\sigmoid(v)/\partial v=1-\sigmoid(v)$ supplies the second line. If the distribution of residual states is held invariant across the paired scenarios, averaging a pointwise nondecreasing conditional mean also gives a nondecreasing marginal mean. This establishes the requested response under the stated model assumptions.

For example, a one-percentage-point increase in the relevant weighted growth input, with magnitude coefficient $0.03$ and no occurrence change, multiplies expected inflow by $e^{0.03}\approx1.03045$. It implies about a 3.05\% rise, not a three-percentage-point change in transaction probability. Units must be fixed before interpreting any estimated coefficient. Applying a coefficient estimated on decimal growth to a percentage-point input would multiply the intended shock by 100.

There are three important qualifications. First, the result is a ceteris paribus GDP response. A joint scenario with higher GDP, unfavorable FX moves, and changing interest rates can have a different total effect. Second, GDP-dependent volatility adds $\tfrac12\partial R_{rr}/\partial g$ to the log-mean derivative and can overturn the sign unless constrained. Third, if bank capture changes with the scenario, monotonic total receipts need not imply monotonic bank-observed receipts. The initial scenario module holds capture fixed or shows a separate capture sensitivity.

\section{Why a positive one-step state loading is insufficient}
An alternative is to put GDP directly in the state equation,
\begin{equation}
x_t=Ax_{t-1}+Bg_t+\eta_t,\qquad y_t=Hx_t+\epsilon_t.
\end{equation}
A GDP impulse at $t$ changes the expected observation $h$ periods later by
\begin{equation}
\frac{\partial\E[y_{t+h}]}{\partial g_t}=HA^hB.
\label{eq:impulseresponse}
\end{equation}
This follows by substituting the state transition repeatedly. Positive entries in $B$ do not ensure positive $HA^hB$: in a scalar stable model with $A=-0.5$, $B=H=1$, the immediate effect is $1$ and the next effect is $-0.5$. A monotone dynamic system could impose nonnegative $H,A,B$, but that also restricts other dynamics. The explicit nonnegative lag kernel in~\eqref{eq:gdpkernel} is a simpler way to make the promised GDP scenario behavior inspectable while allowing residual states to describe other patterns.

Outflow GDP coefficients should initially be unrestricted or separately constrained for a stated economic reason. Even when both expected inflows and outflows rise, net flow can fall if payments respond more strongly. A universal increase in inflows therefore does not imply greater net FX exposure or greater hedge demand. The later utility calculation is needed precisely because those implications do not follow automatically.

\section{Abduction, intervention, and prediction for a client}
Pearl's Definition~4 and equations (27)--(28) formalize a counterfactual by retaining the exogenous conditions of a unit while replacing a structural equation. In the simple structural model $Y=\beta M+U$, observing $(M=m,Y=y)$ implies $U=y-\beta m$. Under $M=m'$,
\begin{equation}
Y(m')=\beta m'+U=y+\beta(m'-m).
\label{eq:unitlinear}
\end{equation}
This expression differs from assigning the same population mean to every client under $m'$: the factual outcome informs the client's disturbance. With noisy measurements or multiple latent states, abduction yields a posterior over $U$ rather than a single value.

For the proposed model, a coherent procedure is:
\begin{enumerate}[leftmargin=*,itemsep=2pt]
\item Infer the posterior of current client states, common states, and relevant past disturbances from observations available at the decision origin.
\item Specify which macro equations or paths change, and which capture, pricing, and behavioral mechanisms remain invariant.
\item Draw future disturbances under the declared structural assumptions and propagate both the factual and alternative systems using paired draws.
\item Compare flow paths, FX exposures, and client utilities, retaining uncertainty from state estimation and the macro response parameters.
\end{enumerate}
Paired random draws reduce simulation noise in scenario differences. They imply a particular cross-world coupling, however. Marginal observational distributions alone do not identify that coupling or individual treatment effects. If the bank has only a predictive conditional model, the product can still deliver useful conditional scenario comparisons while describing the structural interpretation as an additional assumption.

Preference parameters can initially be held fixed while macro scenarios change forecast cash flows. If stress itself changes risk tolerance, financing constraints, offer availability, or the bank's pricing, those mechanisms need to change in the corresponding module. Modular design permits those extensions while making the original assumption visible.

\chapter{From currency exposure to the value of an FX hedge}
\label{ch:utility}
\section{Define the exposure before defining preferences}
The forecasting model produces receipts and payments. A hedge pays against a currency exposure at specified settlement dates. To connect them, define residual native-currency exposure in bucket $h$ as
\begin{equation}
X_{ich}=F^+_{ic,t+h}-F^-_{ic,t+h}+L_{ich}-K_{ich},
\end{equation}
where $L$ is a separately verified balance or commitment allocated to the bucket and $K$ is the already hedged amount. The sign convention is positive for a net foreign-currency receipt. Each balance must be allocated once; adding the entire balance to every future bucket double counts it. Commitments already included in the flow forecast must likewise be reconciled rather than added again. When total $F$ is unobserved, evaluate bank-observed exposure or a declared capture scenario from Chapter~\ref{ch:observation}.

Let $S_{ch}$ denote USD per unit of currency $c$ at settlement. A forward sale of $q_{ch}>0$ units at contractual rate $K^{\rm fwd}_{ch}$ has USD payoff $q_{ch}(K^{\rm fwd}_{ch}-S_{ch})$. A forward purchase corresponds to negative $q$. If $W^0$ is the client's unhedged USD financial outcome, terminal wealth under offer $j$ is
\begin{equation}
W(j)=W^0+\sum_{c,h}q_{j,ch}(K^{\rm fwd}_{j,ch}-S_{ch})-C_j,
\label{eq:wealth}
\end{equation}
with all cash amounts discounted or accumulated to the same valuation date. Costs $C_j$ include fees and any explicitly modeled funding, collateral, or operational costs. Existing hedge payoffs belong in $W^0$ or in the net exposure accounting, consistently. They must not appear twice.

Forward rates are not generally forecasts of future spot rates. In a frictionless illustration with spot $S_0$ in USD per EUR and continuously compounded domestic and foreign rates $r_d,r_f$, borrowing EUR 1, converting to $S_0$ USD, and investing dollars creates $S_0e^{r_dH}$ USD against an obligation $e^{r_fH}$ EUR. A no-arbitrage forward therefore satisfies $K=S_0e^{(r_d-r_f)H}$. Actual executable quotes and client-specific costs should be used in the product model. Expected utility uses physical outcome probabilities; pricing relations and quoted forward points are separate inputs.

\section{Why a corporation may value hedging}
A corporate utility function should have an economic interpretation. If markets are frictionless and shareholders can diversify, reducing firm-level cash-flow variance need not create firm value. \citet{smith1985} analyze reasons why hedging can matter, including taxes, costs of financial distress, and managerial incentives. Their Sections~II--IV provide distinct mechanisms rather than a universal assertion that every firm should minimize variance.

Their tax argument can be understood with state-contingent pretax firm value $V_s$, tax liability $T(V_s)$, and positive state prices $q_s$. Let $Q=\sum_s q_s$ and $\widetilde p_s=q_s/Q$. A zero-cost hedge can, in a complete-market illustration, replace $V_s$ by the constant $\overline V=\sum_s\widetilde p_sV_s$, since
$\sum_s q_s(\overline V-V_s)=0$. If the tax function is convex,
\begin{equation}
T(\overline V)\leq\sum_s\widetilde p_sT(V_s).
\end{equation}
After-tax state-priced value therefore weakly increases under this smoothing, before hedging costs. This is Jensen's inequality applied under normalized state prices, not an assertion about physical expected taxes under an arbitrary hedge. Actual tax schedules, loss carryforwards, and market incompleteness can change the result; historical tax examples in the paper are not current tax guidance.

A convex liquidity-shortfall cost provides another transparent illustration. If a required payment is $D$, let the loss be $\lambda(D-W)_+^2$. The function is convex in wealth shortfall, so a mean-preserving spread of $W$ increases expected loss under the defining convex-order comparison. Reducing variance alone is insufficient to establish this for arbitrary distributions; the full distribution matters. The cost coefficient and payment boundary require client information or sensitivity analysis.

Managerial preferences add a different mechanism. If manager compensation is $w(V)$ and personal utility is $u(w)$, then
\begin{equation}
\frac{d^2u(w(V))}{dV^2}=u''(w)(w')^2+u'(w)w''.
\end{equation}
Concave personal utility gives a negative first term; convex option-like compensation can give a positive second term. The effective preference over firm value need not be concave. This distinction, emphasized in Smith and Stulz's managerial-incentive analysis, argues for estimating or constraining an effective corporate objective rather than assuming a universal risk-aversion coefficient.

\section{Costly external finance generates a concave continuation value}
\label{sec:froot}
\citet{froot1993} connect risk management to investment and financing decisions. The locally stored NBER working paper is the 1992 version of the 1993 journal article; its basic model in Section~3 and FX extension in Section~5 motivate the following derivation. Let $w$ be internal funds available for investment, $I$ investment, and $e=I-w$ external financing. Let $f(I)$ be gross investment payoff with $f''(I)<0$, and $C(e)$ a convex incremental financing cost with $C''(e)>0$. At an interior solution with positive external financing, incremental continuation profit is
\begin{equation}
P(w)=\max_I\{f(I)-I-C(I-w)\}.
\label{eq:frootvalue}
\end{equation}
The cost of investment is $I$; financing friction is the additional $C$. The first-order condition is
\begin{equation}
f'(I^*)-1=C'(I^*-w).
\end{equation}
Differentiating with respect to $w$ gives
\begin{equation}
f''(I^*)\frac{dI^*}{dw}
=C''(e^*)\left(\frac{dI^*}{dw}-1\right),
\qquad
\frac{dI^*}{dw}=\frac{C''}{C''-f''}\in(0,1).
\end{equation}
Additional internal funds support additional investment but also substitute for costly financing. The envelope theorem and another derivative yield
\begin{align}
P'(w)&=C'(e^*),\\
P''(w)&=C''(e^*)\left(\frac{dI^*}{dw}-1\right)
=\frac{f''C''}{C''-f''}<0.
\label{eq:frootcurvature}
\end{align}
Internal funds have diminishing marginal value because shortfalls force costly financing or foregone investment. If a feasible fair hedge changes $w$ without changing its expectation and reduces its risk in convex order, concavity gives a higher $\E[P(w)]$. Costs, hedge-market incompleteness, and changing investment opportunities can prevent full hedging from being optimal.

The source defines $P(w)$ as net investment profit. Total value including the available internal funds is $G(w)=w+P(w)$, so $G'(w)=1+P'(w)$ and $G''(w)=P''(w)$. When a fair hedge preserves expected funds, maximizing $\E[P(w)]$ and $\E[G(w)]$ is equivalent. With fees or changes in expected payoff, the linear funds term must be retained. For an FX payoff $q(K-S)$ and deterministic transaction cost $C_H(q)$, the total-value objective is
\begin{equation}
\max_q\E\left[G\{w^0+q(K-S)-C_H(q),Z\}\right],
\end{equation}
where $Z$ can describe state-dependent investment opportunities. An interior first-order condition is
\begin{equation}
\E\left[(1+P_w\{w(q),Z\})\{K-S-C_H'(q)\}\right]=0.
\label{eq:froothedge}
\end{equation}
The hedge is valuable when it pays in states where internal funds have high marginal value. Minimizing variance of foreign revenues alone need not accomplish that. If foreign-currency payments or attractive foreign investment opportunities rise when receipts rise, the appropriate hedge can be much smaller than gross receipts. The FX extension and Appendix of Froot and coauthors make this investment-opportunity channel explicit. Our one-period notation isolates the mechanism needed here; it does not reproduce all equilibrium cases of their model.

In implementation, estimating $f$ and $C$ for every private company would require far more information than bank transaction histories. The continuation-value derivation supplies an economic reason for concavity. A parsimonious effective utility is a staged approximation to that reason, with sensitivity ranges and later estimation from choices.

\section{A tractable certainty equivalent and its exact domain}
Assume constant absolute risk aversion (CARA),
\begin{equation}
u_i(W)=-e^{-\gamma_iW},\qquad\gamma_i>0,
\end{equation}
where $W$ is measured in USD and $\gamma_i$ has units USD$^{-1}$. The certainty equivalent is the sure wealth giving the same utility:
\begin{equation}
\operatorname{CE}_i(W)=-\frac1{\gamma_i}\log\E[e^{-\gamma_iW}],
\label{eq:ce}
\end{equation}
provided the expectation is finite. If $W\sim\N(\mu,\sigma^2)$, completing the square gives
\begin{align}
-\gamma w-\frac{(w-\mu)^2}{2\sigma^2}
&=-\frac{(w-\mu+\gamma\sigma^2)^2}{2\sigma^2}
-\gamma\mu+\frac{\gamma^2\sigma^2}{2},\\
\E[e^{-\gamma W}]&=e^{-\gamma\mu+\gamma^2\sigma^2/2},\\
\operatorname{CE}(W)&=\mu-\frac{\gamma}{2}\sigma^2.
\label{eq:canormal}
\end{align}
Thus mean--variance utility is exact for normal terminal wealth under CARA. With nonnormal wealth, expanding the cumulant-generating function around zero gives
\begin{equation}
\operatorname{CE}(W)=\kappa_1-\frac\gamma2\kappa_2
+\frac{\gamma^2}{6}\kappa_3-\frac{\gamma^3}{24}\kappa_4+\cdots,
\end{equation}
when this expansion exists. Truncation after the variance is an approximation whose quality depends on risk aversion and higher cumulants.

This distinction is material for the proposed cash-flow model. Stochastic foreign receipts multiplied by stochastic FX rates are generally nonnormal. Moreover, an unbounded lognormal payment can make $\E[e^{-\gamma W}]$ infinite: if $W=-L$ and $L$ is lognormal, the positive exponential moment $\E[e^{\gamma L}]$ diverges. A finite Monte Carlo calculation cannot establish finiteness of that population expectation. The first utility release should therefore use one of two explicitly declared specifications: a mean--variance preference functional with checked finite second moments, or exact expected utility on an economically justified bounded/default-limited wealth model with a finite exponential moment. A finite scenario distribution defines its own finite-support decision model; it is not proof that an unbounded underlying CARA expectation exists.

For finite weighted scenarios, exact CARA certainty equivalent is
\begin{equation}
\widehat{\operatorname{CE}}(W)
=-\gamma^{-1}\log\left(\sum_s w_s e^{-\gamma W_s}\right).
\end{equation}
Evaluate it with log-sum-exp arithmetic to prevent overflow. If $a=\max_s(-\gamma W_s)$, the log term equals $a+\log\sum_sw_se^{-\gamma W_s-a}$. Numerical stability does not remove model risk in extreme scenarios; inspect the states that dominate the exponential weighting.

\section{Deriving hedge size under mean--variance preferences}
Let $R=S-\E[S]$ be the centered FX price change in a single settlement bucket. Write $\delta=K-\E[S]$ and terminal wealth as
\begin{equation}
W(q)=W^0-qR+q\delta-C_H(q).
\end{equation}
Set $v=\V(R)>0$ and $c=\Cov(W^0,R)$. With a deterministic cost,
\begin{align}
\E[W(q)]&=\E[W^0]+q\delta-C_H(q),\\
\V(W(q))&=\V(W^0)+q^2v-2qc.
\end{align}
Relative mean--variance value is consequently
\begin{equation}
\Delta\operatorname{CE}(q)
=q\delta-C_H(q)+\gamma qc-\frac\gamma2q^2v.
\label{eq:hedgegain}
\end{equation}
For quadratic cost $C_H(q)=\kappa q^2/2$, differentiating yields
\begin{equation}
q^*=\frac{\delta+\gamma c}{\gamma v+\kappa}.
\label{eq:hedgeopt}
\end{equation}
This is a unique optimum when the denominator is positive. The term $\gamma c$ represents hedging demand; $\delta$ represents the physical expected carry of the quoted forward. The bank should distinguish those channels and constrain speculative positions according to the product mandate.

With deterministic net receipt $X$ and $W^0=XS+\text{constant}$, $c=Xv$. A fair-in-expectation forward ($\delta=0$) with no costs gives $q^*=X$: the natural hedge is net exposure, not gross receipts. With costs the hedge is smaller. If $q$ is constrained to $[0,X]$, the scalar concave objective is maximized by clipping the unconstrained solution to that interval. A fixed participation cost creates a separate discrete decision: compare optimized hedging value less the fixed cost with zero.

With uncertain receipts, the covariance must be computed jointly. Let $S=\mu_S+s$, $X=\mu_X+x$, with centered $s,x$. Since $W^0=XS$,
\begin{equation}
\Cov(XS,S)=\mu_X\V(S)+\mu_S\Cov(X,S)+\E[xs^2].
\label{eq:randomexposure}
\end{equation}
Expanding $(\mu_X+x)(\mu_S+s)s$ and taking expectations proves the expression. The last term vanishes under centered joint Gaussianity but need not vanish for cash-flow mixtures. Substituting $\E[X]$ into a deterministic exposure formula can therefore misprice risk even if the expected flow is accurate.

For several currency/maturity hedge instruments, let $R$ be a vector of centered settlement prices, $\Sigma=\V(R)$, $c=\Cov(R,W^0)$, and quadratic cost matrix $K_C\succeq0$. The value difference is
\begin{equation}
\Delta\operatorname{CE}(q)=q'\delta+\gamma q'c
-\tfrac12q'(\gamma\Sigma+K_C)q,
\end{equation}
so
\begin{equation}
q^*=(\gamma\Sigma+K_C)^{-1}(\delta+\gamma c)
\end{equation}
when the matrix is positive definite and no constraints bind. With notional, maturity, eligibility, and collateral limits, solve the corresponding constrained concave quadratic program. Singular covariance requires explicit handling of redundant instruments, not an unexplained numerical ridge that changes the recommended position.

\section{A worked exporter example}
Return to the exporter with deterministic EUR 1 million receipts and EUR 0.8 million payments at the same date. Let the future USD/EUR rate have standard deviation $0.05$, set the forward rate equal to the physical expected spot for this illustration, and take $\gamma=10^{-6}$ USD$^{-1}$. The net unhedged position is EUR 200,000, with USD standard deviation $200{,}000\times0.05=10{,}000$ and variance $10^8$. A full net hedge removes that variance, giving mean--variance benefit
\begin{equation}
\frac12\gamma\sigma^2=\tfrac12\times10^{-6}\times10^8=50\ \text{USD}.
\end{equation}
A USD 30 fixed cost leaves USD 20 of value under these illustrative preferences. A USD 80 cost makes this offer unattractive. Hedging EUR 1 million instead leaves a net short EUR 800,000, with USD standard deviation 40,000 and a risk penalty of USD 800. Gross-flow ranking would give a poor recommendation in this example.

The modest USD 50 willingness to pay is a consequence of the illustrative risk coefficient, not a claim about actual corporate willingness to pay. A firm with a binding liquidity covenant or high marginal external-financing cost may value protection much more. The standalone utility module can report sensitivity across effective risk aversion, funding-shortfall costs, capture shares, and existing hedges before enough choice data exist to estimate preferences. This is a useful phase-two analytical product even before a discrete-choice model is fitted.

\chapter{An independent discrete-choice model for FX hedges}
\label{ch:choice}
\section{Utility measures value; choice also reflects frictions}
The utility module can compare a forward sale, a smaller hedge, and no new trade using the same cash-flow scenarios. A choice model adds the fact that clients have heterogeneous preferences and that the bank does not observe every consideration. It should consume scenario-derived values and offer attributes through a stable interface, so that updating a cash-flow model does not require changing the mathematical definition of choice.

At decision occasion $t$, client $i$ chooses from the actual available set $\Avail_{it}$. Include an outside alternative $0$ and define deterministic utility in normalized units as
\begin{equation}
v_{itj}=\lambda\,\Delta\operatorname{CE}_{itj}(\gamma_i)
+\alpha_j+z_i'\Gamma a_{itj}+\theta'a_{itj}-\psi'b_{itj},
\label{eq:choiceutility}
\end{equation}
where $a_{itj}$ describes product and service attributes and $b_{itj}$ describes implementation burdens. The parameter $\lambda>0$ translates dollar certainty equivalent to utility units. Fees already included in the certainty equivalent should not receive a second independent fee penalty without an explicit behavioral reason. Normalize the outside utility to zero after defining all quantities relative to it. If it includes unobserved competitor trades, its economic meaning is a composite alternative.

For a first release, use a small catalog of meaningful currency/maturity/notional offers. If notional is actually a continuous client choice, a coarse catalog approximates that continuous problem. A later model can combine participation choice with conditional continuous sizing. The forecast module remains unchanged in either case.

\section{Deriving McFadden's conditional logit}
\citet[Section I, equations (12)--(13)]{mcfadden1974} derives conditional logit from random utility with independent type-I extreme-value disturbances. Let
\begin{equation}
U_j=v_j+\epsilon_j,\qquad
F(\epsilon)=\exp\{-e^{-\epsilon/\sigma_\epsilon}\},\quad\sigma_\epsilon>0.
\end{equation}
Conditional on $\epsilon_j=e$, alternative $j$ wins when
$\epsilon_k\leq e+v_j-v_k$ for every $k\ne j$. Thus
\begin{align}
P_j&=\int f(e)\prod_{k\ne j}F(e+v_j-v_k)\,de\\
&=\int_{-\infty}^{\infty}\frac1{\sigma_\epsilon}e^{-e/\sigma_\epsilon}
\exp\left\{-e^{-e/\sigma_\epsilon}
 \sum_k e^{(v_k-v_j)/\sigma_\epsilon}\right\}\,de.
\end{align}
Substitute $z=e^{-e/\sigma_\epsilon}$, so $de=-\sigma_\epsilon dz/z$. Reversing the limits gives
\begin{equation}
P_j=\int_0^\infty e^{-zA_j}\,dz=\frac1{A_j}
=\frac{e^{v_j/\sigma_\epsilon}}{\sum_{k\in\Avail}e^{v_k/\sigma_\epsilon}},
\quad A_j=\sum_k e^{(v_k-v_j)/\sigma_\epsilon}.
\label{eq:logit}
\end{equation}
The derivation needs independence and the specified disturbance distribution. Utility differences and a utility scale, rather than absolute utility levels, determine probabilities. Adding the same constant to every $v_j$ changes nothing; multiplying every $v_j$ and $\sigma_\epsilon$ by the same positive number also changes nothing. We therefore normalize $\sigma_\epsilon=1$ and an outside intercept.

For a linear index $v_{itj}=x_{itj}'\theta$, define $d_{itj}=1$ for the chosen alternative and zero otherwise. The log likelihood is
\begin{equation}
\ell(\theta)=\sum_{i,t}\left[
\sum_{j\in\Avail_{it}}d_{itj}x_{itj}'\theta
-\log\sum_{j\in\Avail_{it}}e^{x_{itj}'\theta}\right].
\label{eq:logitlik}
\end{equation}
Differentiation gives the score and Hessian underlying McFadden's Section~II, equations (18)--(20):
\begin{align}
\nabla\ell&=\sum_{i,t}\left(x_{it,y_{it}}-\overline x_{it}\right),
\quad\overline x_{it}=\sum_jP_{itj}x_{itj},\\
\nabla^2\ell&=-\sum_{i,t}\sum_jP_{itj}
(x_{itj}-\overline x_{it})(x_{itj}-\overline x_{it})'.
\label{eq:logithess}
\end{align}
The Hessian is negative semidefinite because each summand is a covariance matrix. Strict identification requires variation in within-choice-set differences; a client attribute identical across alternatives cancels unless interacted with an alternative attribute. Complete separation can lead to unbounded coefficients despite concavity, so finite-solution diagnostics and regularization are necessary. If $\gamma$ enters a nonlinear certainty equivalent, the objective need not remain globally concave in the enlarged parameter vector.

In the simplest two-alternative illustration, let the hedge have normalized utility $0.6$ and the outside option zero. Its predicted share is $e^{0.6}/(1+e^{0.6})\approx0.646$. If a higher cost reduces hedge utility to $0.2$, its predicted share becomes about $0.550$. These calculations show how utility translates into probabilities under the assumptions; they do not identify the effect of changing prices in observational sales records.

\section{Substitution patterns and mixed logit}
Conditional logit implies
\begin{equation}
\frac{P_j}{P_k}=e^{v_j-v_k},
\end{equation}
independent of other available alternatives. This independence of irrelevant alternatives (IIA) can be implausible for several similar forward offers. With equal utilities for ``hedge'' and ``no hedge,'' each gets probability $1/2$. Adding an economically indistinguishable second hedge offer gives each of three alternatives probability $1/3$, so total hedge probability rises to $2/3$. Independent alternative shocks treat duplicate offers as independent attractions. Such a menu change may not reflect real treasury decisions.

\citet[Sections I--III]{mcfadden2000} use mixtures of multinomial logits to allow flexible substitution and heterogeneity. Let $\beta_i=\mu+L\nu_i$, with $\nu_i$ drawn from a specified base density $f(\nu)$. Conditional on $\beta_i$, use the logit formula; then integrate:
\begin{equation}
P_{itj}=\int
\frac{e^{x_{itj}'\beta}}{\sum_{k\in\Avail_{it}}e^{x_{itk}'\beta}}
\,dG(\beta).
\label{eq:mixedlogit}
\end{equation}
Risk aversion and price sensitivity can have sign-constrained parameterizations, such as $\gamma_i=e^{a_i}$, rather than normal coefficients that permit economically unintended signs. Such constraints are hypotheses to validate, and require sufficient variation to estimate.

The mechanism for correlated alternatives is explicit: conditional tastes generate a preference for similar attributes in several offers. Integrating over shared tastes creates dependence in substitution. Theorem~1 of the paper establishes a broader approximation result for suitable random-utility models under compactness, continuity, and absence of ties. Its proof uses finite polynomial approximation and small extreme-value perturbations. For the narrower mechanism used here, a simple argument explains the flexibility. Given random deterministic utilities $V_j(\nu)$ with a unique maximizer almost surely,
\begin{equation}
\frac{e^{V_j(\nu)/\epsilon}}{\sum_ke^{V_k(\nu)/\epsilon}}
\longrightarrow\ind\{j=\argmax_kV_k(\nu)\}
\quad\text{as }\epsilon\downarrow0.
\end{equation}
Dividing numerator and denominator by the largest exponential proves the pointwise limit. Since the probability lies between zero and one, dominated convergence allows integration over $\nu$. This establishes approximation of the corresponding random-utility choice probabilities for a fixed finite menu; it is not a proof of every uniform-approximation condition in the paper's general theorem. Flexibility alone does not establish identifiability or predictive superiority on bank data.

Repeated choices from the same client inform persistent tastes. The panel likelihood must integrate a product:
\begin{equation}
L_i=\int\prod_{t\in\mathcal T_i}
P_{it,y_{it}}(\beta)\,dG(\beta).
\label{eq:panellik}
\end{equation}
Using $\prod_t\int P_{it,y_{it}}(\beta)dG(\beta)$ instead redraws tastes independently each time and changes the model. With $R$ fixed draws $\nu_r$, simulated likelihood uses
\begin{equation}
\widehat L_i=\frac1R\sum_{r=1}^R
\prod_tP_{it,y_{it}}(\mu+L\nu_r).
\end{equation}
Common draws across a client's occasions preserve taste dependence, and fixed draws across optimizer iterations reduce numerical noise. The simulated probability is unbiased under independent sampling, but $\E[\log\widehat L_i]\leq\log L_i$ by Jensen's inequality. Draw-count stability and independent simulation checks are therefore required before trusting estimated preferences or elasticities. The paper's simulated-likelihood methods motivate this estimator; they do not remove finite-simulation bias by assumption.

\section{What the bank must observe to identify preferences}
An observed hedge purchase is not sufficient to reconstruct the choice set. If the relationship manager offered only a three-month forward, the absence of a six-month purchase says little about the client's preference for six months. If a competitor trade is invisible, a no-purchase label mixes several behaviors. The initial estimable target may consequently be ``accept our recorded offer versus all other outcomes,'' conditional on offer eligibility and outcome observation, rather than market-wide hedge choice.

Prices and offers are often endogenous. A salesperson may quote a better spread to a difficult client or contact a client after hearing about a private contract. Let unobserved need be $\zeta$. If quoted price depends on $\zeta$ and $\zeta$ also enters utility, then $\E[\epsilon_j\mid\text{price},x]\ne0$, contradicting the simple exogeneity assumption. Adding price to a logit does not identify a causal price response. Possible remedies are randomized variation within approved quoting practices, valid excluded cost shifters with a defensible structural treatment, or a joint offer/pricing model with explicitly justified assumptions. A statistical instrument needs relevance and exclusion; merely having a macro variable does not make it valid.

There is also scale confounding. A dollar-valued certainty equivalent is translated by $\lambda$, while $\gamma$ controls the variance penalty. In a narrow dataset where every offer changes mean and variance proportionally, only combinations such as $\lambda\gamma$ may be well estimated. Variation in costs, maturity, notional, and exposure is needed to distinguish them. Initial releases can fix a small declared grid of economically interpretable risk coefficients and report sensitivity, then estimate a parsimonious preference distribution once the choice records support it. A million clients with no declined-offer data do not resolve this problem.

The first choice likelihood treats scenario-derived values as inputs available at the offer date. Those values are estimates of what the bank believes about exposure, not necessarily the client's private beliefs. If the client has a large unobserved contract, the choice residual reflects both private information and preference. Separating risk aversion from belief differences requires additional data or assumptions. Parameter uncertainty from the flow model should be propagated by repeated fitted-model draws or another declared approximation. This statistical propagation preserves modular software and estimation stages; it does not force joint training of cash flows and utility.

\section{Product counterfactuals can be delivered separately}
Once preference and pricing assumptions are credible, the choice module can compare the same forecast distribution under a different fee, maturity, hedge fraction, or available menu. A price counterfactual recalculates client wealth and the utility index, then recomputes choice probabilities. It can be evaluated without refitting the time-series model if the price change does not affect the underlying operating cash flows during the horizon.

That separability is an assumption of timing. A hedge could alter collateral usage or financing capacity, which in turn changes investment and later cash flows. The first product treats operating forecasts as predetermined for a short decision horizon and includes known funding effects in utility. A longer-horizon structural model can feed financing consequences back into the state equations. Such feedback would be a later research extension, with additional data requirements, rather than a hidden implication of a one-period logit.

\chapter{Sales allocation and the incremental value of contact}
\label{ch:outreach}
\section{A likely purchase and a useful sales call are different targets}
The utility and choice modules can identify clients likely to value an FX hedge. The bank still needs to know where a call changes the outcome. A client who routinely hedges through an automated channel may have high purchase probability but low incremental benefit from relationship-manager time. Another client with moderate exposure may act only after receiving an explanation. Contact propensity must therefore be evaluated as a treatment rather than substituted for product preference.

Let $X$ contain information available before outreach, $A$ indicate assignment to outreach, and $R$ be realized net contribution over a fixed horizon. Net contribution should include spreads or fees less relevant execution, credit, funding, service, and contact costs. Define $\mu_a(x)=\E[R\mid A=a,X=x]$ and $e(x)=\Pr(A=1\mid X=x)$. Under consistency, conditional exchangeability $(R(0),R(1))\perp A\mid X$, overlap $0<e(X)<1$, and no relevant interference, the conditional effect is
\begin{equation}
\tau(x)=\mu_1(x)-\mu_0(x).
\end{equation}
These assumptions are demanding for historical discretionary sales activity. A randomized pilot creates more credible evidence than increasingly elaborate adjustment for undocumented salesperson judgment.

\section{A doubly robust score and its derivation}
\citet{wager2021policy} develops policy learning from observational data using orthogonal scores. Their equation (4) and policy-learning sections motivate the augmented inverse-propensity score
\begin{equation}
\Gamma=\widehat\mu_1(X)-\widehat\mu_0(X)
+\frac{A}{\widehat e(X)}\{R-\widehat\mu_1(X)\}
-\frac{1-A}{1-\widehat e(X)}\{R-\widehat\mu_0(X)\}.
\label{eq:aipw}
\end{equation}
To see the robustness, condition on $X=x$. If the propensity is correct, then
\begin{align}
\E\!\left[\frac A{e(x)}(R-\widehat\mu_1)\mid x\right]
&=\mu_1-\widehat\mu_1,\\
\E\!\left[\frac{1-A}{1-e(x)}(R-\widehat\mu_0)\mid x\right]
&=\mu_0-\widehat\mu_0.
\end{align}
Substituting into~\eqref{eq:aipw} leaves $\mu_1-\mu_0$. If instead both outcome regressions are correct, each residual has zero conditional mean in its treatment group, so the augmentation terms vanish even with a misspecified propensity. With nuisance limits $\widetilde\mu_a,\widetilde e$, the conditional bias is exactly
\begin{equation}
\E[\widetilde\Gamma\mid x]-\tau(x)
=(\widetilde e-e)
\left\{\frac{\widetilde\mu_1-\mu_1}{\widetilde e}
+\frac{\widetilde\mu_0-\mu_0}{1-\widetilde e}\right\}.
\end{equation}
The product form explains why modest nuisance errors can be less damaging than in a direct plug-in estimate. It does not correct unmeasured confounding or absent overlap. Cross-fitting estimates nuisance functions on other folds so that each client's score is evaluated without using its own outcome to fit those functions. Time-respecting splits are additionally needed for a genuinely prospective policy evaluation.

For a binary policy $\pi(x)$, incremental value relative to never contacting is
\begin{equation}
V(\pi)-V(0)=\E[\pi(X)\tau(X)],
\qquad
\widehat{V(\pi)-V(0)}=\frac1n\sum_i\pi(X_i)\widehat\Gamma_i.
\end{equation}
The source paper's objective written using $(2\pi-1)\Gamma$ is an affine transformation of this one and selects the same unconstrained policy class. Capacity and eligibility constraints restrict that class in the banking implementation. Learning and evaluating the policy on the same outcomes biases the reported value upward; a held-out evaluation or randomized deployment is required.

\section{Capacity, product offers, and phased use}
If every contact uses equal time and effects are additive, the solution to~\eqref{eq:capacityintro} selects the largest positive $\tau_i$ up to capacity. With unequal time costs, the binary problem is a knapsack problem; sorting by $\tau_i/d_i$ is exact for a divisible relaxation but not generally for indivisible calls. Territory ownership, language, expertise, cooldown periods, and client eligibility can be expressed as assignment constraints.

The forecasting and choice models can supply useful features for $\tau$: expected unhedged exposure, uncertainty, estimated client value, acceptance probability, and timing. They do not mathematically determine the causal effect of contact. A heuristic interim queue can rank eligible clients by estimated client benefit and expected bank contribution, with its status stated clearly. The later pilot tests whether that queue generates incremental net value compared with current coverage and simpler rules.

Randomize outreach assignment among eligible clients within practical strata such as salesperson, industry, and predicted opportunity. Define the outcome horizon and contact cost before starting, preserve assignment even if contact is not completed, and use intention-to-treat as the primary assignment effect. Repeated contacts need a cooldown or a sequential design. Cluster assignment or inference when clients share a corporate group or when salesperson capacity creates interference. An experimental sample-size calculation must use pilot outcome variance, economically meaningful minimum effects, and the actual clustering; no credible sample size can be supplied from client count alone.

Unobserved competitor trades affect interpretation here as well. The bank can measure incremental contribution at its own institution if outcome capture is reliable. It cannot infer changes in the client's total market hedging from its own trade data. Both can be worthwhile business questions, but they require different outcome data.

\chapter{Validation, delivery phases, and commercial feasibility}
\label{ch:delivery}
\section{What would count as success?}
The proposal is useful if it produces better financial decisions with the information the bank actually possesses. Low average prediction error is one piece of that evidence. A hedge application also needs calibrated uncertainty, credible covariance, honest exposure scope, and useful behavior under scenarios. A sales product ultimately needs incremental contribution at realistic capacity. These criteria should be fixed before selecting a model on the final holdout period.

The main forecast comparison should use out-of-time predictive distributions of currency-specific receipts and payments, cumulative net exposure, and relevant liquidity shortfalls. The first business pilot should use incremental net contribution and client-value safeguards at the intended outreach budget. Predictive accuracy cannot serve as a substitute for that randomized commercial result.

The proposed architecture addresses the problem through explicit mechanisms. Pooling reduces the number of independently estimated client parameters. State dynamics distinguish persistent activity from payment noise. A hurdle likelihood accommodates zero months and positive skew. Joint covariance preserves natural hedges. Macro response parameters permit transparent scenario changes. The independent utility module values offers from financial outcomes, and the choice module accounts for heterogeneity and available alternatives. These are reasons the design can represent the required questions. Whether it forecasts or sells well is an empirical question for the following evaluation.

\section{Chronological evaluation with 36 dates}
Use rolling forecast origins so that every fitted parameter, transform, peer statistic, and feature uses only information available at its origin. A feasible illustrative split is months 1--18 for initial fitting, origins 18--24 for development subject to label availability, and origins 25--30 for final six-month evaluation through month 36. For a six-month target, a development origin of 24 uses outcomes through month 30; those outcomes must not be used when pretending to deploy a model at month 25. One clean implementation therefore freezes hyperparameter selection before the final test origin, using only development targets already matured at that time. Another reserves a later deployment date and documents the changed test window. The precise schedule should be constructed from label-availability dates, not just row partitions.

For clarity, a conservative final six-month test can freeze all selection using origins no later than 18 and outcomes through month 24, then evaluate origins 25--30 with predeclared rolling refitting rules. Shorter horizons can provide more origins, but overlapping targets create dependence. This small number of common test dates gives limited evidence about macro regimes. Client-level bootstrap intervals alone can become misleadingly narrow because all clients share the same calendar shocks. Report paired differences by origin and industry, use time-block or two-way dependence-aware uncertainty where justified, and state when the number of time blocks is too small for reliable asymptotics.

Operational forecast tests must use macroeconomic information or macro forecasts available at the forecast origin. Conditioning on realized future GDP or FX is an oracle diagnostic of the conditional response model, not a deployable forecast result. Report it separately. Scenario-conditioned performance and the performance of an unconditional operational forecast answer different questions.

Two generalization questions should be evaluated separately: forecasts for existing clients with prior history, and forecasts for new clients or newly active currencies. Hold out complete entities or corporate groups for the second question. Do not let related subsidiaries appear in both sides of a group-level holdout when their shared information would make the exercise unrealistic.

\section{A constructed set of practical competitors}
The initial model must face inexpensive rules that a treasury practitioner could actually use. The following competitors are defined operationally, and each should produce a comparable predictive distribution using training-only residual or pooled uncertainty estimates where needed:
\begin{enumerate}[leftmargin=*,itemsep=3pt]
\item \textbf{Last observed amount:} forecast the most recent complete month's inflow and outflow in each currency. It is a strong persistence rule for stable clients.
\item \textbf{Trailing mean:} forecast each coordinate by its last six complete months' mean, with the window treated as a development choice. It resists a single unusual month.
\item \textbf{Seasonal naive:} use the same calendar month one year earlier when available, with a pooled fallback for shorter histories. It tests whether annual scheduling alone explains the signal.
\item \textbf{Recency--frequency--amount rule:} estimate occurrence from recent positive-month frequency and positive size from recent nonzero amounts, shrunk toward the industry/currency mean. It is a transparent intermittent-flow adversary.
\item \textbf{Pooled lagged regression:} use industry, currency, lagged gross flows, and calendar variables with regularization. It tests the incremental value of latent states.
\item \textbf{Transparent Gaussian state-space model:} fit a small shared level/seasonal model with explicit covariance. It tests whether the hurdle and richer hierarchy add useful information.
\end{enumerate}
The plain proposed model is the hierarchical hurdle state-space model without optional neural parameter generation. Enhanced candidates add residual common factors, richer macro response, or a neural generator only in controlled comparisons. DeepAR and deep state-space models are flexible literature-based challengers, with their covariate information and parameter-search budgets matched fairly to the structured model.

Compare each method conditionally for stable versus volatile clients, dense versus intermittent flows, long versus short relationships, exporters versus importers, natural hedges versus one-sided currency positions, missing observations, large clients, and macro stress or structural-change periods represented in the data. An average gain cannot conceal a failure on a commercially material segment. Underperformance against a simple rule in such a segment blocks promotion for that segment until repaired or routed to an explicit fallback. The competitor set is a falsification check, not a test set to repeatedly optimize against.

\section{Metrics matched to the financial task}
For marginal distribution accuracy, use log predictive density and the continuous ranked probability score (CRPS). If $Z,Z'$ are independent draws from forecast $F$, then
\begin{equation}
\operatorname{CRPS}(F,y)=\E|Z-y|-\tfrac12\E|Z-Z'|.
\end{equation}
The first term rewards proximity to the realization; the second adjusts for the forecast's own spread. For a sample of $M$ paths it can be estimated using pairwise distances or a sorted-sample algorithm. The hurdle distribution has a point mass at zero and a continuous positive density; log-score evaluation must respect that mixed measure rather than adding an arbitrary density epsilon.

For intervals, report coverage and width by currency, horizon, industry, and activity regime. For joint forecasts, evaluate realized distributions of $Y^+-Y^-$, cumulative exposure, and hedge P\&L, as well as joint scoring rules. One example is the energy score,
\begin{equation}
\operatorname{ES}(F,\bm y)=\E\|\bm Z-\bm y\|-\tfrac12\E\|\bm Z-\bm Z'\|.
\end{equation}
Currency and horizon scales must be declared before applying a vector norm. Because a joint score can be insensitive to some dependence errors in high dimensions, direct covariance and exposure-tail diagnostics remain necessary. For a decision quantile $q_\alpha$, use pinball loss
\begin{equation}
L_\alpha(y,q_\alpha)=(\alpha-\ind\{y<q_\alpha\})(y-q_\alpha).
\end{equation}
Expected pinball loss is minimized at an $\alpha$-quantile: where differentiable, its derivative with respect to $q$ is $F_Y(q)-\alpha$. At atoms, the subgradient condition gives the usual quantile interval. This supplies a direct interpretation for one-sided funding or receipt risk forecasts.

For choice, evaluate held-out log likelihood, probability calibration, and substitution under observed menu variation. Choice accuracy is a weak criterion when one alternative dominates. For utility calculations, inspect realized hedge costs and P\&L, but do not treat one favorable currency move as evidence that a hedge was ex ante optimal. The objective concerns the distribution known at the decision date. For outreach, evaluate net contribution under randomized assignment, with capacity and contact costs included.

Scenario validation is different from forecast validation. Test the imposed GDP monotonicity over all supported horizons, check units and lags, inspect unrestricted-response sensitivity, and test whether the specification extrapolates sensibly. Passing these checks establishes that the implementation obeys the proposed assumptions. It does not empirically validate an unobserved intervention. A causal interpretation requires separate identification evidence.

\section{Phases that deliver value independently}
\subsection{Phase 0: reliable observations and a coverage baseline}
Deliver a dated native-currency inflow/outflow panel with entity boundaries, reconciliation checks, zero/missing distinctions, and observation coverage. Construct the six simple competitors, and begin logging complete offer sets, quotes, declines, and contact assignment. Resolve whether the earliest product speaks about bank-observed activity or supported total-wallet estimates. Exit this phase when a sampled set of entities reconciles to source systems and historical features can be reproduced at their actual availability dates.

\subsection{Phase 1: cash-flow forecasts and macro scenarios}
Deliver joint receipt/payment paths, forecast intervals, exposure summaries visible to the bank, and a scenario control for GDP and other macro inputs. Start with the small structured model and the explicit GDP response kernel. Report predictive versus imposed versus identified scenario interpretations. This phase has no dependency on estimating utility or choice. Promotion requires out-of-time calibration and usefulness against the constructed competitors, including material client segments; an unvalidated neural enhancement remains optional.

\subsection{Phase 2: FX exposure and client-value calculations}
Add verified balances, commitments, existing hedges, indicative product quotes, and settlement-bucket accounting. Deliver comparative hedge values under declared preferences and capture sensitivities. Use the analytic hedge formula as a reference case and scenario evaluation for richer exposures. Show which uncertainty comes from flows, FX, capture, or preferences. This phase can support an informed treasury discussion before behavior is estimated. Promotion requires accurate cash-flow-to-payoff accounting and economically coherent sensitivity, not a claim that assumed preferences are revealed preferences.

\subsection{Phase 3: estimated choice and controlled outreach}
Fit a parsimonious conditional logit once actual choice sets and dated prices exist. Add mixed logit only when repeated choices and heterogeneity warrant it. Design approved variation or a defensible identification strategy for price effects. Run an outreach experiment to compare the proposed queue with current coverage and simpler rules. Promotion requires uncertainty-supported incremental net value, appropriate product eligibility, and credible measurement of outcomes. A well-calibrated forecast alone cannot satisfy this criterion.

\section{A production design proportionate to the project}
Use monthly batch estimation and scoring initially, with event-triggered refreshes only when newer underlying data exist. Share group parameters, store compact client states, and simulate detailed paths for eligible currency-active clients or shortlisted opportunities. A million clients with 1,000 paths, 6 months, 10 currencies, 2 directions, and 4-byte values would require 480 billion bytes, roughly 480 GB, before exchange rates and metadata. Streaming summaries and reproducible scenario seeds reduce storage; detailed path retention should be selective and justified by downstream use.

Each module should accept versioned inputs and return units, timestamps, assumptions, uncertainty summaries, and reason codes understandable to a relationship manager. A useful client view would say that observed euro receipts and supplier payments imply an uncertain net exposure in a specified period, show the effect of a selected GDP scenario, and identify the information still needed to verify a hedge amount. It should not expose model-fitting jargon in place of the financial explanation.

Monitor observation coverage, entity-link changes, new-currency activity, forecast calibration, dependence errors, and choice-set recording. Track sales policy changes because they alter future labels and observed bank capture. A rise in the bank's share after outreach can look like an increase in underlying client cash flows unless capture and treatment are distinguished. Model refreshes need an as-of evaluation against the previous version and the same practical baselines.

\section{What remains uncertain, and the next justified decision}
The proposal is mathematically capable of generating the requested flow scenarios and independent product choices. Its most consequential empirical uncertainties are total-wallet visibility, the stability and identification of macro responses, tail dependence between flows and FX, and the availability of complete choice records. These should determine the first data work, rather than selecting a large architecture before seeing the observations.

No empirical superiority, return on investment, or causal GDP elasticity is reported here. The immediate justified action is a bounded data audit and retrospective forecasting prototype under a predeclared evaluation plan. Exact compute budgets, training searches, random-seed replication, and promotion thresholds should be fixed after the panel dimensions, active currencies, and data quality are known. That avoids treating convenient numerical settings or calendar estimates as scientific defaults.

\chapter{Source scope and the contribution of the proposal}
\label{ch:sources}
\section{How the twelve works support the design}
The survey concentrates on methods that change the proposed model or its interpretation. It is a selective methodological synthesis rather than an exhaustive history of cash-flow forecasting, macroeconometrics, corporate finance, and marketing. Each cited work has a local full-text copy, and the main text supplies the derivations needed for the proposed use. Claims about bank performance, causal GDP responses, and client preferences remain to be tested with the bank's data.

\citet{montero2021global} supplies the local-versus-global forecasting distinction and the reason to ask how much structure can be shared. Chapter~\ref{ch:dynamics} explains the scope of its expressiveness result and derives a concrete Gaussian pooling rule. \citet{kalman1960} supplies recursive linear estimation; the Gaussian specialization, multistep forecast, and likelihood are derived there. \citet{shumway1982} supplies the smoothing--EM connection; the sufficient statistics and transition/noise updates are reconstructed explicitly. Their direct relevance is learning persistent client conditions without fitting a free forecasting model to only 36 observations.

\citet{stock2002} supplies factor compression for many predictors, with asymptotic conditions that prevent interpreting client count as macroeconomic history. \citet{salinas2020deepar} supplies shared autoregressive distributional forecasting and a useful neural competitor. \citet{rangapuram2018} supplies a shared parameter generator combined with analytic state filtering. Chapter~\ref{ch:dynamics} gives their relevant likelihood constructions and distinguishes the source models from the proposed multivariate hurdle extension. Chapter~\ref{ch:joint} adds the bank-specific observation model, dependence structure, and approximate inference.

\citet{pearl2009} provides the language and structural definitions required to distinguish observational conditioning, intervention, and unit counterfactuals. Chapter~\ref{ch:cf} derives the confounding example, adjustment argument, and abduction procedure. The nonnegative GDP lag kernel and its monotonicity proof are local design choices, not conclusions of that source.

\citet{smith1985} explains tax, distress, and managerial channels for corporate hedging. \citet{froot1993} explains the value of internal funds when investment and financing interact. Chapter~\ref{ch:utility} reconstructs the relevant curvature argument and then derives a separate tractable utility model, clearly stating when normality or finite-moment conditions are needed. The source papers motivate economic preferences; they do not identify this bank's clients' risk coefficients.

\citet{mcfadden1974} supplies the random-utility derivation of conditional logit and its likelihood geometry. \citet{mcfadden2000} supplies mixed-logit integration, panel heterogeneity, and simulation methods. Chapter~\ref{ch:choice} derives the formulas used, illustrates substitution limitations, and explains what observed offers and prices must identify. Finally, \citet{wager2021policy} supplies the doubly robust policy-learning approach in Chapter~\ref{ch:outreach}. It is relevant to the eventual use of sales capacity, independently of whether cash flows and product choice are predicted accurately.

\section{Versions, local access, and limits of coverage}
All twelve PDFs are stored under \path{dcos/papers/salespitch/}, with filenames of the form \texttt{Title FirstAuthor(year).pdf}. Colons in titles are rendered as spaced hyphens for portable filenames. The publication year in a filename identifies the bibliographic work; the accompanying \texttt{manifest-v2.json} records the exact downloaded version, source URL, checksum, and technical anchors. Earlier proposal files are preserved.

Several version differences deserve explicit notice. The DeepAR bibliography describes the four-author 2020 journal article; the downloaded arXiv version 3 dated 2019 lists Salinas, Flunkert, and Januschowski. The exposition relies on the inspected autoregressive and likelihood construction in that version. The Froot--Scharfstein--Stein technical source is NBER Working Paper 4084 from May 1992, preceding the 1993 journal publication. The deep state-space PDF is the official supplementary long version, which includes the non-Gaussian inference discussion absent from the short conference paper. The Kalman copy is a readable transcription of the published article. Scanned originals for Shumway--Stoffer, Smith--Stulz, Froot and coauthors, and McFadden's chapter have local OCR sidecars used for search; equations were reconstructed mathematically rather than copied from imperfect OCR.

The literature search included backward references and bounded forward-citation queries. Citation metadata are dated 8 October 2026; some metadata requests were rate-limited. Citation counts and venue indicators are recorded when available in the review files and are not treated as evidence that a method is correct. The limited forward samples do not establish exhaustive coverage or a global ranking of influential follow-up work. No retraction notice was encountered in the inspected primary records; this is not an exhaustive registry clearance.

Important neighboring areas include structural macroeconomic shock identification, richer corporate financing models, probabilistic intermittent-demand forecasting, nested choice structures, joint continuous/discrete hedge sizing, and recent general-purpose time-series models. They remain meaningful competitors or extensions. The current proposal deliberately does not claim to identify a policy shock from 36 GDP observations, optimize an entire corporate balance sheet, or establish that a foundation model improves client-level hedge decisions. The separate omission register explains the resulting coverage risks and the evidence that would justify extending the survey.

\section{What is new here, and what has been checked}
The contribution is an integrated, phased design for the bank's information constraints: native-currency gross flows; an explicit partial-observation boundary; a pooled dynamic hurdle distribution; an inspectable macro response channel; a separate corporate-utility calculation; available-set choice; and eventual causal sales allocation. The architecture allows a useful forecast and scenario product to be delivered before preferences, competitor choices, or outreach effects are identified.

Companion deterministic checks verify selected mathematical identities, including Gaussian filtering and smoothing, GDP response signs, net-exposure covariance, hedge optimization, logit derivatives, the panel-mixture distinction, and the doubly robust score. These checks support the displayed derivations. They are not bank-data experiments or a trained implementation of the full proposal. The document has been compiled and its rendered pages inspected; reader comprehension and empirical product acceptance remain separate reviews.

\bibliographystyle{plainnat}
\bibliography{references}
\end{document}
